%2multibyte Version: 5.50.0.2960 CodePage: 65001
%%\newenvironment{proof}[1][Proof]{\noindent\textbf{#1.} }{\ \rule{0.5em}{0.5em}}


\documentclass[12pt,notitlepage,a4paper]{article}
%%%%%%%%%%%%%%%%%%%%%%%%%%%%%%%%%%%%%%%%%%%%%%%%%%%%%%%%%%%%%%%%%%%%%%%%%%%%%%%%%%%%%%%%%%%%%%%%%%%%%%%%%%%%%%%%%%%%%%%%%%%%%%%%%%%%%%%%%%%%%%%%%%%%%%%%%%%%%%%%%%%%%%%%%%%%%%%%%%%%%%%%%%%%%%%%%%%%%%%%%%%%%%%%%%%%%%%%%%%%%%%%%%%%%%%%%%%%%%%%%%%%%%%%%%%%
\usepackage{amsmath}
\usepackage{amssymb}
\usepackage{amsfonts}
\usepackage[onehalfspacing,nodisplayskipstretch]{setspace}
\usepackage[a4paper,nohead,left=1in,right=1in,top=1in,bottom=1in]{geometry}

\setcounter{MaxMatrixCols}{10}
%TCIDATA{OutputFilter=LATEX.DLL}
%TCIDATA{Version=5.50.0.2960}
%TCIDATA{Codepage=65001}
%TCIDATA{<META NAME="SaveForMode" CONTENT="1">}
%TCIDATA{BibliographyScheme=Manual}
%TCIDATA{Created=Monday, December 20, 2010 11:42:45}
%TCIDATA{LastRevised=Wednesday, April 08, 2026 16:54:23}
%TCIDATA{<META NAME="GraphicsSave" CONTENT="32">}
%TCIDATA{<META NAME="DocumentShell" CONTENT="Standard LaTeX\Standard LaTeX Article">}
%TCIDATA{Language=American English}
%TCIDATA{CSTFile=LaTeX article (bright).cst}
%TCIDATA{ComputeDefs=
%$\lambda $
%$a^{\Lambda }$
%$a_{t}^{\Lambda }$
%$\nu (\{n_{0},n_{0}$
%$\alpha $
%}


\newtheorem{theorem}{Theorem}
\newtheorem{acknowledgement}[theorem]{Acknowledgement}
\newtheorem{algorithm}[theorem]{Algorithm}
\newtheorem{axiom}[theorem]{Axiom}
\newtheorem{case}[theorem]{Case}
\newtheorem{claim}[theorem]{Claim}
\newtheorem{conclusion}[theorem]{Conclusion}
\newtheorem{condition}[theorem]{Condition}
\newtheorem{conjecture}[theorem]{Conjecture}
\newtheorem{corollary}[theorem]{Corollary}
\newtheorem{criterion}[theorem]{Criterion}
\newtheorem{definition}[theorem]{Definition}
\newtheorem{example}[theorem]{Example}
\let\oldexample\example
\renewcommand{\example}{\oldexample\normalfont}
\newtheorem{exercise}[theorem]{Exercise}
\newtheorem{lemma}[theorem]{Lemma}
\newtheorem{notation}[theorem]{Notation}
\newtheorem{problem}[theorem]{Problem}
\newtheorem{proposition}[theorem]{Proposition}
\newtheorem{remark}[theorem]{Remark}
\newtheorem{solution}[theorem]{Solution}
\newtheorem{summary}[theorem]{Summary}
\newenvironment{proof}[1][Proof]{\noindent\textbf{#1.} }{\hfill $\square$}
\newcounter{savedtheorem}
% tcilatex.tex - stub for Scientific WorkPlace compatibility

\begin{document}

\title{Proper Calibeating\thanks{%
Previous versions: December 2025. We thank for useful comments and
suggestions. A presentation is available at \texttt{%
http://www.ma.huji.ac.il/hart/pres.html\#calib-beat-p}}}
\author{Dean P. Foster\thanks{%
Department of Statistics, Wharton, University of Pennsylvania, Philadelphia,
and Amazon, New York. \emph{e-mail}: \texttt{dean@foster.net} \ \emph{web
page}: \texttt{http://deanfoster.net}} \and Sergiu Hart\thanks{%
Institute of Mathematics, Department of Economics, and Federmann Center for
the Study of Rationality, The Hebrew University of Jerusalem. \emph{e-mail}: 
\texttt{hart@huji.ac.il} \ \emph{web page}: \texttt{%
http://www.ma.huji.ac.il/hart}}}
\maketitle

\begin{abstract}
\end{abstract}

\tableofcontents

%TCIMACRO{%
%\TeXButton{References Without Numbers}{\def\@biblabel#1{#1\hfill}
%\def\thebibliography#1{\section*{References}
%\addcontentsline{toc}{section}{References}
%\list
%{}{
%\labelwidth 0pt
%\leftmargin 1.8em
%\itemindent -1.8em
%\usecounter{enumi}}
%\def\newblock{\hskip .11em plus .33em minus .07em}
%\sloppy\clubpenalty4000\widowpenalty4000
%\sfcode`\.=1000\relax\def\baselinestretch{1}\large \normalsize}
%\let\endthebibliography=\endlist} }%
%BeginExpansion
\def\@biblabel#1{#1\hfill}
\def\thebibliography#1{\section*{References}
\addcontentsline{toc}{section}{References}
\list
{}{
\labelwidth 0pt
\leftmargin 1.8em
\itemindent -1.8em
\usecounter{enumi}}
\def\newblock{\hskip .11em plus .33em minus .07em}
\sloppy\clubpenalty4000\widowpenalty4000
\sfcode`\.=1000\relax\def\baselinestretch{1}\large \normalsize}
\let\endthebibliography=\endlist
%EndExpansion
%TCIMACRO{%
%\TeXButton{define shfrac}{\newcommand{\shfrac}[2]{\ensuremath{{}^{#1} \hspace{-0.04in}/_{\hspace{-0.03in}#2}}}} }%
%BeginExpansion
\newcommand{\shfrac}[2]{\ensuremath{{}^{#1} \hspace{-0.04in}/_{\hspace{-0.03in}#2}}}
%EndExpansion
%TCIMACRO{%
%\TeXButton{def \T \B}{\newcommand\T{\rule{0pt}{2.6ex}}
%\newcommand\B{\rule[-1.2ex]{0pt}{0pt}}}}%
%BeginExpansion
\newcommand\T{\rule{0pt}{2.6ex}}
\newcommand\B{\rule[-1.2ex]{0pt}{0pt}}%
%EndExpansion

\section{Introduction\label{s:intro}}

Forecasters---whether of weather, elections, or sporting events---issue
probabilistic predictions that are then compared with the outcomes that
materialize. How should one evaluate such a forecaster? A natural approach
is to use a \emph{scoring rule}: a function that assigns a loss to each
forecast--outcome pair. A scoring rule is \emph{proper} if the forecaster
minimizes expected loss by reporting truthfully; the standard choice is the
\emph{quadratic} (Brier 1950) scoring rule (see Savage 1971, Schervish
1989, and the comprehensive treatment in Gneiting and Raftery 2007).

A classic result of DeGroot and Fienberg (1983) is that the Brier score
decomposes as\footnote{%
For forecasts that are sorted into bins by their announced values; see
Section \ref{s-s:sequences} for the precise setup.}%
\begin{equation}
\mathcal{B}=\mathcal{K}+\mathcal{R}\text{,}  \label{eq:BKR-intro}
\end{equation}%
where $\mathcal{K}$ is the \emph{calibration} score---measuring how close
the forecasts are to the realized frequencies---and $\mathcal{R}$ is the
\emph{refinement} score---measuring how well the forecaster sorts outcomes
into informative bins. A surprising result of Foster and Vohra (1998) is
that one can generate forecasts that are \emph{guaranteed} to be calibrated
(i.e., $\mathcal{K}$ converges to zero), no matter what the outcomes turn
out to be (see also Foster 1999, Foster and Hart 2021, and the survey of
Olszewski 2015).

In Foster and Hart (2023), we introduced the notion of
\textquotedblleft calibeating\textquotedblright : beating a reference
forecaster at its own game by achieving a Brier score no worse than its
refinement score---thus gaining calibration without losing expertise. We
showed that calibeating can be achieved by a simple deterministic online
procedure, by a stochastic procedure that is itself calibrated, and by
deterministic continuously calibrated procedures. All these results,
however, were stated exclusively for the Brier score.

But there are infinitely many proper scoring rules, and it would be
disturbing if a forecaster's guarantees evaporated merely because the
evaluator chose a different proper loss. This raises the question:%
\begin{quote}
\emph{If a procedure is calibrated (or calibeats) under the Brier score, do
these guarantees persist under every other proper scoring rule?}
\end{quote}
This paper answers that question for proper Lipschitz scoring rules. The
answer is positive for calibration and negative for calibeating.

We will call a procedure \emph{proper-calibrated} (respectively, \emph{%
proper-calibeating}) if the corresponding guarantee holds simultaneously
for all proper Lipschitz scoring rules, after normalizing by the Lipschitz
constant (see Section \ref{s:proper} for the precise definitions). We
restrict attention to Lipschitz scoring rules; this excludes rules, such as
the logarithmic scoring rule, that may take the value $+\infty $ on the
boundary of the simplex. The restriction is natural for our purpose, because
we seek guarantees that hold uniformly over a class of scoring rules.

\subsection{Calibration transfers for free\label{s-s:intro-calibration}}

We show, first, that calibration transfers for free: every $\varepsilon $%
-calibrated procedure is automatically $\varepsilon $-\emph{proper}%
-calibrated. The key observation is that for any proper $M$-Lipschitz
scoring rule $L$,%
\begin{equation}
\mathcal{K}_{t}^{L}(\mathbf{c};\mathbf{f})\leq M\,\mathcal{K}_{t}(\mathbf{%
c};\mathbf{f})\text{;}  \label{eq:KL-K-intro}
\end{equation}%
that is, the $L$-calibration score is bounded by the Lipschitz constant
times the quadratic calibration score. In particular, all known calibrated
forecasting procedures---the stochastic procedures of Foster and Vohra
(1998), the approachability-based constructions of Blackwell (1956), Foster
(1999), Kakade and Foster (2008), Hart (2022), and Mannor and Stoltz (2010),
and the deterministic continuously calibrated procedures of Foster and Hart
(2021)---are already proper-calibrated, with no modification needed.

\subsection{Calibeating does not transfer\label{s-s:intro-not-transfer}}

Our second result is perhaps a surprise: unlike calibration, calibeating
does \emph{not} transfer. We exhibit a concrete example (Section \ref%
{s-s:example}) in which a forecasting sequence $\mathbf{c}$ calibeats a
reference $\mathbf{b}$ under the quadratic scoring rule, yet \emph{fails}
to calibeat $\mathbf{b}$ under the $2$-spherical scoring rule---which is
also proper and Lipschitz. Both sequences are perfectly calibrated; the
failure is purely in the refinement comparison.

The reason for the asymmetry is instructive. Calibration involves only $%
\mathcal{K}^{L}$, which is controlled by $\mathcal{K}$ via the bound (\ref%
{eq:KL-K-intro}). Calibeating involves $\mathcal{B}^{L}-\mathcal{R}^{L}$,
and there is no analogous inequality relating $\mathcal{R}^{L}$ to $%
\mathcal{R}$: refinement scores under different proper rules are not
comparable in general.

\subsection{Proper calibeating\label{s-s:intro-proper-calibeating}}

We then provide three positive results on proper calibeating (i.e.,
calibeating simultaneously under all proper Lipschitz scoring rules).

First, the simple deterministic procedure of Theorem 3 of Foster and Hart
(2023)---which forecasts the running average action in the reference
forecaster's bin---turns out to be proper-calibeating. The proof relies on
a new bound showing that online refinement $\widetilde{\mathcal{R}}^{L}$
converges to offline refinement $\mathcal{R}^{L}$ at rate $O((\ln t)/t)$
for every proper $M$-Lipschitz $L$, generalizing the quadratic-case result
of Foster and Hart (2023, Proposition 2).

Second, a stochastic procedure that calibeats the \emph{joint} binning $%
\mathbf{b}\times \mathbf{c}$ is both $(\delta ,B)$-proper-calibeating and $%
\delta $-proper-calibrated. This identifies the right structural condition:
it is not enough to calibeat the reference forecaster and to be calibrated;
what matters is calibeating the \emph{joint} binning, which always refines
the forecast partition and thereby restores the refinement structure needed
for the argument. The proof chains the decomposition $\mathcal{B}^{L}=%
\mathcal{K}^{L}(\mathbf{c};\mathbf{b}\times \mathbf{c})+\mathcal{R}^{L}(%
\mathbf{b}\times \mathbf{c})$ with the bound (\ref{eq:KL-K-intro}).

Third, there exists a deterministic continuously proper-calibrated procedure
that proper-calibeats. Since fractional binnings do not refine the
forecast-binning, the exact decomposition $\mathcal{B}^{L}=\mathcal{K}%
^{L}+\mathcal{R}^{L}$ breaks down. We prove an approximate version:%
\begin{equation*}
\left\vert \mathcal{B}_{t}^{L}(\mathbf{c})-\left( \mathcal{K}_{t}^{L}(%
\mathbf{c};\mathbf{f})+\mathcal{R}_{t}^{L}(\mathbf{f})\right) \right\vert
<2M\delta
\end{equation*}%
for any $\delta $-local fractional binning under an $M$-Lipschitz rule.
This generalizes the quadratic approximate decomposition from the errata
(Foster and Hart 2026), and is the key technical tool for this result.

All three results extend to simultaneously proper-calibeating $N$ reference
forecasters, by applying the above to the product binning $\mathbf{b}%
^{1}\times \cdots \times \mathbf{b}^{N}$ (cf.\ Lee, Noarov, Pai, and Roth
2022, who independently study multicalibeating via online multiobjective
optimization).

\subsection{Related work\label{s-s:intro-related}}

Several concurrent lines of work address---from different angles---the
question of universality across scoring rules.

The most directly related is that of Chen, Huang, Jordan, and Luo (2026),
who reduce calibeating to regret minimization and obtain calibeating rates
for general proper losses (including Brier, log, and mixable losses). Their
approach works loss by loss: for each proper loss $L$, they build a
procedure achieving $\mathcal{B}^{L}\leq \mathcal{R}^{L}+O((\log T)/T)$.
Our approach is different: we show that a \emph{single} procedure (the
existing Brier-score procedures) achieves proper calibeating \emph{%
simultaneously} for all proper $1$-Lipschitz rules $L$. The uniformity
is a consequence of the bound (\ref{eq:KL-K-intro}), which has no analogue
for refinement. They also handle the logarithmic rule, which is not
Lipschitz and so lies outside our class $\mathcal{L}_{1}$.

Kleinberg, Leme, Schneider, and Teng (2023) introduce
\textquotedblleft U-calibration,\textquotedblright\ asking what forecasting
criterion guarantees low regret for an unknown agent who may use any bounded
scoring rule. They show that calibration suffices but is not necessary, and
achieve $O(\sqrt{T})$ U-calibration error. Our work is complementary: we
are not introducing a new criterion but showing that existing
calibration and calibeating procedures already provide uniform guarantees
over all proper Lipschitz rules. Gopalan, Kalai, Reingold, Sharan, and
Wieder (2022) show that multicalibration implies good performance for all
convex Lipschitz loss functions (\textquotedblleft
omniprediction\textquotedblright ); our universality is analogous but
concerns \emph{evaluation rules} (scoring rules), not decision-making
losses.

Dimitriadis, Gneiting, and Jordan (2023) decompose the continuous ranked
probability score into miscalibration, discrimination, and uncertainty
components; Popordanoska, Uria, and Blundell (2023) show that every proper
score decomposes into a proper calibration error and a refinement term via a
Bregman divergence. Both study the calibration--refinement decomposition for
general proper rules from a statistical perspective; our work provides the
game-theoretic counterpart.

Blasiok, Gopalan, Hu, and Nakkiran (2023) prove a converse direction: local
optimality of a proper loss implies smooth calibration. Luo, Senapati, and
Sharan (2025) achieve simultaneous swap regret minimization for all proper
losses with smooth univariate form via KL-calibration, providing yet another
route to universality.

\subsection{Outline of the paper\label{s-s:intro-outline}}

Section \ref{s:setup} sets up the framework: scoring rules, divergences,
and the Brier, calibration, and refinement scores. Section \ref{s:proper}
introduces the \textquotedblleft proper\textquotedblright\ concepts
(proper-calibration and proper-calibeating). Section \ref{s:proper
calibration} proves that calibration is automatically proper-calibration.
Section \ref{s:proper calibeating} presents the counterexample showing that
calibeating does not imply proper calibeating, establishes the three routes
to proper calibeating, and extends the results to multicalibeating. The
appendices collect the background on scoring rules and the analysis of the
joint-binning condition.

\section{The Setup\label{s:setup}}

The setup is as in our previous papers, Foster and Hart (2018, 2021, 2023,
2026), except that the set of forecasts $C$ is now specified as the
appropriate probability simplex $\Delta $.

Let $A$ be a finite set of \emph{actions}, and let $\Delta :=\{c\in \mathbb{R%
}_{+}^{A}:\sum_{a\in A}c(a)=1\}$, the simplex of probability distributions
over the set $A$, be the set of \emph{forecasts.} We identify the elements
of $A$ with the unit vectors of $\Delta $.

\subsection{Scoring Rules\label{s-s:scoring rules}}

A \emph{scoring rule} $L_{A}:A\times \Delta \rightarrow \mathbb{R}$
associates to every forecast $c\in \Delta $ and every realized action $a\in
A $ a loss\footnote{%
There are scoring rules, such as the logarithmic scoring rule, that allow
the loss to be infinite; as they are not Lipschitz-continuous, we will not
deal with them here (for \textquotedblleft
log-calibeating,\textquotedblright\ see Appendix A.9 in [FH2022]).} $%
L_{A}(a,c)$. The function $L_{A}$ is linearly extended to $L:\Delta \times
\Delta \rightarrow \mathbb{R}$ by $L(d,c):=\sum_{a\in A}d(a)L_{A}(a,c)$;
thus, $L(d,c)=\mathbb{E}_{a\sim d}\left[ L_{A}(a,c)\right] $ is the expected
loss when the realization is distributed according to\footnote{%
For $a$ in A we thus have $L_{A}(a,c)=L(\mathbf{1}_{a},c).$} $d\in \Delta .$
Writing $\mathbf{L}(c)$ for the vector $(L_{A}(a,c))_{a\in A}$ in $\mathbb{R}%
^{A}$, we thus have%
\begin{equation}
L(d,c)=d\cdot \mathbf{L}(c)  \label{eq:d*Lambda(c)}
\end{equation}%
for every $c,d\in \Delta .$

A scoring rule $L$ is \emph{proper} if $L(d,c)\geq L(d,d)$ for every $c,d\in
\Delta $; i.e., forecasting the correct distribution minimizes the expected
loss (if $L(d,c)>L(d,d)$ for all $c\neq d$, then $L$ is \emph{strictly proper%
}). To measure the deviation from the perfect forecast $c=d$, define the $L$%
\emph{-divergence}\footnote{%
At times this is written $D(d||c)$ (as is the case for the Kullback-Leibler
divergence).} $D\equiv D^{L}:\Delta \times \Delta \rightarrow \mathbb{R}$ by%
\begin{equation}
D(d,c)\,:=\,L(d,c)-L(d,d).  \label{eq:D-le}
\end{equation}%
Thus, $L$ is proper if and only if $D$ is always $\geq 0.$ The standard
scoring rule is the \emph{quadratic (Brier) scoring rule}, whose divergence
is $D(d,c)=\left\Vert c-d\right\Vert ^{2}$ (and $L_{A}(a,c)=-2c(a)+\left%
\Vert c\right\Vert ^{2}$).

A scoring rule $L$ is $M$\emph{-Lipschitz} for some finite $M$ if%
\begin{equation*}
\left\Vert \mathbf{L}(c)-\mathbf{L}(c^{\prime })\right\Vert \leq M\left\Vert
c-c^{\prime }\right\Vert
\end{equation*}%
for all $c,c^{\prime }\in \Delta $, where $\left\Vert \cdot \right\Vert $
denotes the standard Euclidean norm.

We list several useful properties (see Appendix \ref{s-a:scoring rules} for
details):

\begin{proposition}
\label{p:L}(i) If $L$ is a proper scoring rule then the function $H\equiv
H^{L}:\Delta \rightarrow \mathbb{R}$ (\textquotedblleft the $L$%
-entropy\textquotedblright ) given by $H(c):=L(c,c)$ is concave, and we have 
\begin{eqnarray*}
L(d,c) &=&H(c)-(c-d)\cdot \mathbf{L}(c)\;\;\;\text{and} \\
D(d,c) &=&H(c)-H(d)-(c-d)\cdot \mathbf{L}(c)
\end{eqnarray*}%
for every $c,d$ in $\Delta $.

(ii) If $L$ is a proper $M$-Lipschitz scoring rule then%
\begin{eqnarray}
0 &\leq &D(d,c)\leq M\left\Vert c-d\right\Vert ^{2}\;\;\;\text{and}
\label{eq:L} \\
\left\vert D(d,c)-D(d,c^{\prime })\right\vert &=&\left\vert
L(d,c)-L(d,c^{\prime })\right\vert \leq M\left\Vert c-c^{\prime }\right\Vert
\label{eq:D-D-L}
\end{eqnarray}%
for every $c,c^{\prime },d$ in $\Delta $.
\end{proposition}

\subsection{Scores for Sequences of Forecasts\label{s-s:sequences}}

Let $t$ be the horizon; for $s=1,...,t$, let $a_{s}\in A$ be the action, $%
c_{s}\in \Delta $ the forecast, and $i_{s}\in I$ the \textquotedblleft
bin\textquotedblright\ (for some set of bins $I$).\footnote{%
We abstract away from the way the binning is determined. The
\textquotedblleft standard\textquotedblright\ binning, by forecast, is given
by $\mathbf{i=c}$, i.e., $i_{s}=c_{s}$ for every $s$.} As in FH2023, we
write $\mathbf{a}^{t}$ for $(a_{s})_{s=1}^{t}$ and $\mathbf{a}$ for $%
(a_{s})_{s=1}^{\infty },$ and similarly for the other sequences.\footnote{%
While usually only $\mathbf{a}^{t},\mathbf{c}^{t},...$ are needed, for
convenience we may often use $\mathbf{a},\mathbf{c},...$ $.$} Assume that
the binning is a refinement of the standard binning that is generated by the
forecasts; i.e., all forecasts in the same bin $i$ have the same $c$
(formally: $i_{s}=i_{r}$ implies $c_{s}=c_{r}$).\footnote{%
Formally, the forecast is \emph{measurable} with respect to the binning.} A
scoring rule $L$, with corresponding divergence function $D^{L}$, generates
the following scores, which we refer to as $L$\emph{-Brier}, $L$\emph{%
-calibration}, and $L$\emph{-refinement}:\footnote{%
The $\mathcal{B}$-score depends on the sequence $\mathbf{c},$ the $\mathcal{R%
}$-score on the sequence $\mathbf{i}$, and the $\mathcal{K}$-score on both
(of course, all of them depend on the sequence $\mathbf{a}$ as well). When $L
$ is the standard quadratic score we drop the superscript $L$.\newline
Since we use different forecasting sequences and different binning
sequences, precise and explicit notation is needed (so as to avoid errors,
as in Foster and Hart 2023; see 2026 Errata).}%
\begin{eqnarray*}
\mathcal{B}_{t}^{L}\equiv \mathcal{B}_{t}^{L}(\mathbf{c}) &%
%TCIMACRO{\TeXButton{:=}{{\;:=\;}}}%
%BeginExpansion
{\;:=\;}%
%EndExpansion
&\frac{1}{t}\sum_{s=1}^{t}D^{L}(a_{s},c_{s}) \\
\mathcal{K}_{t}^{L}\equiv \mathcal{K}_{t}^{L}(\mathbf{c};\mathbf{i}) &%
%TCIMACRO{\TeXButton{:=}{{\;:=\;}}}%
%BeginExpansion
{\;:=\;}%
%EndExpansion
&\frac{1}{t}\sum_{s=1}^{t}D^{L}(\bar{a}_{t}(i_{s}),c_{s}) \\
\mathcal{R}_{t}^{L}\equiv \mathcal{R}_{t}^{L}(\mathbf{i}) &%
%TCIMACRO{\TeXButton{:=}{{\;:=\;}}}%
%BeginExpansion
{\;:=\;}%
%EndExpansion
&\frac{1}{t}\sum_{s=1}^{t}D^{L}(a_{s},\bar{a}_{t}(i_{s})),
\end{eqnarray*}%
where for each $i$ in $I$%
\begin{equation*}
n_{t}(i)\,%
%TCIMACRO{\TeXButton{:=}{{\;:=\;}}}%
%BeginExpansion
{\;:=\;}%
%EndExpansion
\,\left\vert \{s\leq t:i_{s}=i\}\right\vert 
\end{equation*}%
is the number of entries in bin $i$, and%
\begin{equation*}
\bar{a}_{t}(i)\,%
%TCIMACRO{\TeXButton{:=}{{\;:=\;}}}%
%BeginExpansion
{\;:=\;}%
%EndExpansion
\,\frac{1}{n_{t}(i)}\sum_{s\leq t:i_{s}=i}a_{s}
\end{equation*}%
is the average action in bin $i$. For the standard binning given by the
forecasts, i.e., when $\mathbf{i}=\mathbf{c}$, we shorten $\mathcal{K}^{L}(%
\mathbf{c};\mathbf{c})$ to $\mathcal{K}^{L}(\mathbf{c}).$

Let%
\begin{equation*}
\mathcal{H}_{t}^{L}%
%TCIMACRO{\TeXButton{:=}{{\;:=\;}}}%
%BeginExpansion
{\;:=\;}%
%EndExpansion
\frac{1}{t}\sum_{s=1}^{t}H^{L}(a_{s})=\frac{1}{t}\sum_{s=1}^{t}L(a_{s},a_{s})
\end{equation*}%
be the average $L$-entropy, then%
\begin{equation*}
\mathcal{B}_{t}^{L}(\mathbf{c})=\frac{1}{t}\sum_{s=1}^{t}L(a_{s},c_{s})-%
\mathcal{H}_{t}^{L}
\end{equation*}%
and, summing by bins and using the linearity of $L$ in its first argument,%
\begin{eqnarray}
\mathcal{R}_{t}^{L}(\mathbf{i}) &=&\frac{1}{t}\sum_{s=1}^{t}L(a_{s},\bar{a}%
_{t}(i_{s}))-\mathcal{H}_{t}^{L}  \notag \\
&=&\sum_{i\in I}\left( \frac{n_{t}(i)}{t}\right) \left( \frac{1}{n_{t}(i)}%
\sum_{s\leq t:i_{s}=i}L(a_{s},\bar{a}_{t}(i))\right) -\mathcal{H}_{t}^{L} 
\notag \\
&=&\sum_{i\in I}\left( \frac{n_{t}(i)}{t}\right) L(\bar{a}_{t}(i),\bar{a}%
_{t}(i))-\mathcal{H}_{t}^{L}  \notag \\
&=&\sum_{i\in I}\left( \frac{n_{t}(i)}{t}\right) H^{L}(\bar{a}_{t}(i))-%
\mathcal{H}_{t}^{L}.  \label{eq:R-H-H}
\end{eqnarray}

The following is a classic decomposition result (XX):

\begin{theorem}
\label{th:S-decompose}Let the binning sequence $\mathbf{i}$ be a refinement
of the forecasting sequence $\mathbf{c}$; then 
\begin{equation*}
\mathcal{B}_{t}^{L}(\mathbf{c})=\mathcal{K}_{t}^{L}(\mathbf{c};\mathbf{i})+%
\mathcal{R}_{t}^{L}(\mathbf{i}).
\end{equation*}
\end{theorem}

\begin{proof}
Let $c^{i}$ denote the forecast in bin $i$. Summing by bins and again using
the linearity of $L$ in its first argument gives 
\begin{eqnarray*}
\mathcal{B}_{t}^{L}(\mathbf{c}) &=&\sum_{i\in I}\left( \frac{n_{t}(i)}{t}%
\right) \left( \frac{1}{n_{t}(i)}\sum_{s\leq t:i_{s}=i}L(a_{s},c^{i})\right)
-\mathcal{H}_{t}^{L} \\
&=&\sum_{i\in I}\left( \frac{n_{t}(i)}{t}\right) L(\bar{a}_{t}(i),c^{i})-%
\mathcal{H}_{t}^{L},
\end{eqnarray*}%
and%
\begin{eqnarray*}
\mathcal{K}_{t}^{L}(\mathbf{c};\mathbf{i}) &=&\sum_{i}\left( \frac{n_{t}(i)}{%
t}\right) \left[ L(\bar{a}_{t}(i),c^{i})-L(\bar{a}_{t}(i),\bar{a}_{t}(i))%
\right] \\
&=&\sum_{i\in I}\left( \frac{n_{t}(i)}{t}\right) L(\bar{a}%
_{t}(i),c^{i})-\sum_{i\in I}\left( \frac{n_{t}(i)}{t}\right) H^{L}(\bar{a}%
_{t}(i)).
\end{eqnarray*}%
Subtracting gives $\mathcal{B}-\mathcal{K}=\mathcal{R}$ by (\ref{eq:R-H-H}).
\end{proof}

\subsubsection{General Fractional Binning Sequences\label{s-s-s:fractional}}

Following Foster and Hart (2021), we now consider general fractional
binnings. Let $I$ be a finite or countably infinite set of bins;\footnote{%
For pure binnings the number of bins is always finite (up to time $t$, it is
at most $t$).} a \emph{general (fractional) binning sequence} $\mathbf{f}%
=(f_{s})_{s=1,2,...}$specifies in each period $s$ the fraction $f_{s}(i)\geq
0$ that is assigned to each bin $i\in I$, where $\sum_{i\in I}f_{s}(i)=1$;
thus, $f_{s}$ may be viewed as a probability distribution on $I$, i.e., $%
f_{s}\in \Delta (I)$. For example, given a fractional binning $\Pi
=(w_{i})_{i\in I}$ (where $w_{i}:\Delta \rightarrow \lbrack 0,1]$ and $%
\sum_{i\in I}w_{i}(c)=1$ for every $c\in \Delta $; see Foster and Hart 2021)
we put\footnote{%
The resulting binning sequence is thus time-independent: the fraction $%
f_{s}(i)$ depends only on the forecast $c_{s}$ and not on the period $s$.} $%
f_{s}(i)=w_{i}(c_{s})$. When each $f_{s}$ is a unit vector, i.e., in each
period there is a single bin only, we will call the binning sequence \emph{%
pure}.

The definitions of calibration and refinement naturally extend to fractional
binnings:%
\begin{eqnarray*}
\mathcal{K}_{t}^{L}(\mathbf{c};\mathbf{f}) &%
%TCIMACRO{\TeXButton{:=}{{\;:=\;}}}%
%BeginExpansion
{\;:=\;}%
%EndExpansion
&\frac{1}{t}\sum_{i\in I}n_{t}(i)D^{L}(\bar{a}_{t}(i),\bar{c}_{t}(i)) \\
\mathcal{R}_{t}^{L}(\mathbf{f}) &%
%TCIMACRO{\TeXButton{:=}{{\;:=\;}}}%
%BeginExpansion
{\;:=\;}%
%EndExpansion
&\frac{1}{t}\sum_{i\in I}\sum_{s=1}^{t}f_{s}(i)D^{L}(a_{s},\bar{a}_{t}(i)),
\end{eqnarray*}%
where for each $i$ in $I$%
\begin{equation*}
n_{t}(i)\,%
%TCIMACRO{\TeXButton{:=}{{\;:=\;}}}%
%BeginExpansion
{\;:=\;}%
%EndExpansion
\,\sum_{s=1}^{t}f_{s}(i)
\end{equation*}%
is the total weight in bin $i$, and%
\begin{eqnarray*}
\bar{a}_{t}(i) &%
%TCIMACRO{\TeXButton{:=}{{\;:=\;}}}%
%BeginExpansion
{\;:=\;}%
%EndExpansion
&\sum_{s=1}^{t}\left( \frac{f_{s}(i)}{n_{t}(i)}\right) a_{s} \\
\bar{c}_{t}(i) &%
%TCIMACRO{\TeXButton{:=}{{\;:=\;}}}%
%BeginExpansion
{\;:=\;}%
%EndExpansion
&\sum_{s=1}^{t}\left( \frac{f_{s}(i)}{n_{t}(i)}\right) c_{s}
\end{eqnarray*}%
are the average action and the average forecast in bin $i$.

These definitions clearly reduce to the previous definitions when the
binning sequence is pure and it refines the standard by-forecast-binning,
because then each bin contains a single value of forecast. For general
binning sequences, where a bin may contain multiple forecast values, the
decomposition of Theorem \ref{th:S-decompose} no longer holds.\footnote{%
For the quadratic scoring rule there is another such a decomposition, but
with a different definition of refinement, namely, as the average
bin-variance of the \emph{differences} $a_{t}-c_{t}$. Carrying this out for
a proper scoring rule would yield here an average of \emph{differences
between divergences} (which is in general not a divergence). We sidestep
this by using Theorem \ref{th:decompose-delta} below (cf. FH2026).} However,
we will now show that it continues to approximately hold for
\textquotedblleft local\textquotedblright\ binning sequences where the
forecasts in each bin are close to one another (and the scoring rule is
Lipschitz).

Given a (continuous) fractional binning $\Pi =(w_{i})_{i\in I},$ let $\Pi (%
\mathbf{c})$ denote the fractional binning sequence $\mathbf{f}$ where $%
f_{s}=(w_{i}(c_{s}))_{i\in I}\in \Delta (I)$ for every $s$; i.e., in each
period $s$ the fractions are given by the weight functions $w_{i}$ applied
to the forecast $c_{s}.$

Let $\delta >0$; a general binning sequence $\mathbf{f}$ is $\delta $\emph{%
-local with respect to the sequence }$\mathbf{c}$ if for each $i$ there is
an open ball $B(y^{i};\delta )$ with center $y^{i}\in \Delta $ and radius $%
\delta $ such that $f_{s}(i)>0$ implies $\left\Vert c_{s}-y^{i}\right\Vert
<\delta $; i.e., all forecasts in bin $i$ lie in $B(y^{i};\delta ).$ Pure
binning sequences that are a refinement of the forecast-binning are thus $%
\delta $-local for every $\delta >0$. The generalization of the
decomposition theorem \ref{th:S-decompose} is:\footnote{%
This is the counterpart of the Lemma in the Errata for the quadratic scoring
rule: $\mathcal{R}^{\#}=\mathcal{B}-\mathcal{K}$, and so $\mathcal{B}-\left[ 
\mathcal{K}+\mathcal{R}\right] =\mathcal{R}^{\#}-\mathcal{R}$.}

\begin{theorem}
\label{th:decompose-delta}Let $\delta >0$; if the general binning sequence $%
\mathbf{f}$ is $\delta $-local with respect to the forecasting sequence $%
\mathbf{c}$, then%
\begin{equation*}
\left\vert \mathcal{B}_{t}^{L}(\mathbf{c})-\left( \mathcal{K}_{t}^{L}(%
\mathbf{c};\mathbf{f})+\mathcal{R}_{t}^{L}(\mathbf{f})\right) \right\vert
<2M\delta
\end{equation*}%
for every $M$-Lipschitz scoring rule $L$.
\end{theorem}

\begin{proof}
Replacing all $c_{s}$ and $\bar{c}_{t}(i)$ in bin $i$ by $y^{i}$ yields
scores\footnote{%
The refinement score is not affected since it does not depend on the
forecasts. For simplicity we drop the subscript $t$ from $n_{t},\bar{a}_{t},$
and $\bar{c}_{t}$.}%
\begin{eqnarray*}
\mathcal{\hat{B}}_{t}^{L}(\mathbf{c}) &%
%TCIMACRO{\TeXButton{:=}{{\;:=\;}}}%
%BeginExpansion
{\;:=\;}%
%EndExpansion
&\sum_{i\in I}\left( \frac{n(i)}{t}\right) \sum_{s=1}^{t}\left( \frac{%
f_{s}(i)}{n(i)}\right) D(a_{s},y^{i})\text{\ \ \ and} \\
\mathcal{\hat{K}}_{t}^{L}(\mathbf{c};\mathbf{f}) &%
%TCIMACRO{\TeXButton{:=}{{\;:=\;}}}%
%BeginExpansion
{\;:=\;}%
%EndExpansion
&\sum_{i\in I}\left( \frac{n(i)}{t}\right) D(\bar{a}(i),y^{i}).
\end{eqnarray*}%
The proof consists of showing, first, that $\mathcal{\hat{B}}$ is close to $%
\mathcal{B}$ and $\mathcal{\hat{K}}$ is close to $\mathcal{K}$, and second,
that\footnote{%
Since positive fractions of $c_{s}$ may be allocated to more than one bin
and thus replaced by \emph{different} $y^{i}$-s, there is no single
replacement of the forecasting sequence $\mathbf{c}$ that would give (ii) by
the decomposition of Theorem \ref{th:S-decompose}; the proof of (ii) is just
a slight variant.} $\mathcal{\hat{B}-\hat{K}=R}$.

(i) 
\begin{eqnarray*}
\left\vert \mathcal{B}_{t}^{L}(\mathbf{c})-\mathcal{\hat{B}}_{t}^{L}(\mathbf{%
c})\right\vert &<&M\delta \text{\ \ \ and} \\
\left\vert \mathcal{K}_{t}^{L}(\mathbf{c};\mathbf{f})-\mathcal{\hat{K}}%
_{t}^{L}(\mathbf{c};\mathbf{f})\right\vert &<&M\delta .
\end{eqnarray*}%
Indeed, all $c_{s}$ in bin $i$ (i.e., with $f_{s}(i)>0$) satisfy $\left\Vert
c_{s}-y^{i}\right\Vert <\delta $, and thus their average $\bar{c}(i)\equiv 
\bar{c}_{t}(i)$ satisfies $\left\Vert \bar{c}(i)-y^{i}\right\Vert <\delta $
as well. Therefore $\left\vert D(a_{s},c_{s})-D(a_{s},y^{i})\right\vert
<M\delta $ and $\left\vert D(\bar{a}(i),\bar{c}(i))-D(\bar{a}%
(i),y^{i})\right\vert <M\delta $ by (\ref{eq:D-D-L}). Averaging the former
over $s$ and $i$ yields the inequality for $\mathcal{B}$, and averaging the
latter over $i$ yields the inequality for $\mathcal{K}$.

(ii)%
\begin{equation*}
\mathcal{\hat{B}}_{t}^{L}(\mathbf{c})=\mathcal{\hat{K}}_{t}^{L}(\mathbf{c};%
\mathbf{f})+\mathcal{R}_{t}^{L}(\mathbf{f}).
\end{equation*}%
Indeed, we have%
\begin{eqnarray*}
\mathcal{\hat{B}}_{t}^{L}(\mathbf{c}) &=&\sum_{i\in I}\left( \frac{n(i)}{t}%
\right) \sum_{s=1}^{t}\left( \frac{f_{s}(i)}{n(i)}\right) L(a_{s},y^{i})-%
\mathcal{H}_{t}^{L} \\
&=&\sum_{i\in I}\left( \frac{n(i)}{t}\right) L(\bar{a}(i),y^{i})-\mathcal{H}%
_{t}^{L},
\end{eqnarray*}%
and%
\begin{eqnarray*}
\mathcal{\hat{K}}_{t}^{L}(\mathbf{c};\mathbf{f}) &=&\sum_{i\in I}\left( 
\frac{n(i)}{t}\right) \left[ L(\bar{a}(i),y^{i})-L(\bar{a}(i),\bar{a}(i))%
\right] \\
&=&\sum_{i\in I}\left( \frac{n(i)}{t}\right) L(\bar{a}(i),y^{i})-\sum_{i\in
I}\left( \frac{n(i)}{t}\right) H(\bar{a}(i)).
\end{eqnarray*}%
Subtracting gives (ii) by (\ref{eq:R-H-H}).

Combining (i) and (ii) yields the result.
\end{proof}

\section{\textquotedblleft Proper\textquotedblright\ Concepts\label{s:proper}%
}

We will say that a procedure is \textquotedblleft proper\textquotedblright
-calibrated if it is calibrated with respect to every proper Lipschitz
scoring rule $L$; that is, the $L$-calibration score converges to zero as
the horizon increases. Moreover, we require uniformity in the scoring rule $L
$.  Since multiplying a scoring rule by a constant multiplies the
calibration score by that same constant, to get uniform convergence we need
to normalize the scoring functions; for Lipschitz scoring rules, a
convenient normalization is to rescale them by dividing by the Lipschitz
constant, which yields $1$-Lipschitz scoring rules.\footnote{%
Functions that are $1$-Lipschitz are called \textquotedblleft nonexpansive."}
Proper-calibeating will be similarly defined. 

Let thus $\mathcal{L}$ denote the class of all proper Lipschitz scoring
rules, and $\mathcal{L}_{1}$ the subclass of $1$-Lipschitz such rules. We
will say that a procedure is \emph{(uniformly) proper-calibrated /
proper-calibeating} if the corresponding requirement holds simultaneously
for all scoring rules in $\mathcal{L}_{1}$ (we will usually drop the term
\textquotedblleft uniform\textquotedblright ). The formal definitions are as
follows, for $\varepsilon \geq 0$ (when $\varepsilon =0$ we say
\textquotedblleft proper-calibrated/calibeating\textquotedblright\ instead
of \textquotedblleft $0$-proper-calibrated/calibeating\textquotedblright ),
procedure $\sigma $ (deterministic or stochastic) with forecasting sequence $%
\mathbf{c}$:

\begin{itemize}
\item The procedure $\sigma $ is \emph{(uniformly) }$\varepsilon $-\emph{%
proper-calibrated} (following Foster and Vohra 1998) if\footnote{%
The expectation $\mathbb{E}$ is taken over the random forecasts of $\sigma $.%
}%
\begin{equation*}
\varlimsup_{t\rightarrow \infty }\left( \sup_{L\in \mathcal{L}_{1}}\sup_{%
\mathbf{a}_{t}\in A^{t}}\mathbb{E}\left[ \mathcal{K}_{t}^{L}(\mathbf{c})%
\right] \right) \leq \varepsilon ^{2}.
\end{equation*}

\item The procedure $\sigma $ is \emph{(uniformly) continuously }$%
\varepsilon $\emph{-proper-calibrated} (following Foster and Hart 2021) if%
\begin{equation*}
\varlimsup_{t\rightarrow \infty }\left( \sup_{L\in \mathcal{L}_{1}}\sup_{%
\mathbf{a}_{t}\in A^{t}}\mathbb{E}\left[ \mathcal{K}_{t}^{L}(\mathbf{c};\Pi (%
\mathbf{c}))\right] \right) \leq \varepsilon ^{2}
\end{equation*}%
for every continuous binning\footnote{%
In Foster and Hart (2021) it is proved that it suffices for the condition to
hold for one specific continuous binning $\Pi _{0}$.} $\Pi $.

\item Let $B$ be a finite set; the\ $\mathbf{b}$-based procedure $\sigma $
is \emph{(uniformly) }$(\varepsilon ,B)$-\emph{proper-calibeating }%
(following Foster and Hart 2023) if%
\begin{equation*}
\varlimsup_{t\rightarrow \infty }\left( \sup_{L\in \mathcal{L}_{1}}\sup_{%
\mathbf{a}_{t}\in A^{t},\mathbf{b}_{t}\in B^{t}}\mathbb{E}\left[ \mathcal{B}%
_{t}^{L}(\mathbf{c})-\mathcal{R}_{t}^{L}(\mathbf{b})\right] \right) \leq
\varepsilon ^{2}.
\end{equation*}
\end{itemize}

From now on, the plain terms \textquotedblleft
calibration\textquotedblright\ and \textquotedblleft
calibeating\textquotedblright\ refer to the standard quadratic scoring rule
only. 

\bigskip 

\noindent \textbf{Remark. }While the quadratic scoring rule needs rescaling
in order to be included in $\mathcal{L}_{1}$ (because its Lipschitz constant
is $2\sqrt{2}$), its divergence $D(d,c)=\left\Vert d-c\right\Vert ^{2}$
satisfies the crucial inequality (\ref{eq:L}) without rescaling. Therefore
it may be convenient to artificially add the quadratic scoring rule without
rescaling it to $\mathcal{L}_{1}$, so that $\varepsilon $%
-proper-calibration/calibeating includes $\varepsilon $%
-calibration/calibeating (with the same $\varepsilon $).

\section{Proper Calibration\label{s:proper calibration}}

Calibration turns out to always yield proper calibration.

\begin{theorem}
\label{th:proper-calibration}Every procedure that is $\varepsilon $%
-calibrated is $\varepsilon $-\emph{proper}-calibrated, and every procedure
that is continuously $\varepsilon $-calibrated is continuously $\varepsilon $%
-\emph{proper}-calibrated.
\end{theorem}

\bigskip 

This is an immediate consequence of the following:

\begin{proposition}
\label{p:K-KS}Let $L$ be a proper $M$-Lipschitz scoring rule. Then\footnote{%
We do not use a superscript $L$ when $L$ is the standard quadratic scoring
rule.}%
\begin{equation*}
\mathcal{K}_{t}^{L}(\mathbf{c};\mathbf{f})\leq M\,\mathcal{K}_{t}(\mathbf{c};%
\mathbf{f}).
\end{equation*}
\end{proposition}

\begin{proof}
Using Proposition \ref{p:L} (ii) yields 
\begin{eqnarray}
\mathcal{K}_{t}^{L}(\mathbf{c};\mathbf{f}) &=&\frac{1}{t}\sum_{i\in I}\left( 
\frac{n_{t}(i)}{t}\right) D^{L}(\bar{a}_{t}(i),\bar{c}_{t}(i))
\label{eq:KS vs K} \\
&\leq &\frac{1}{t}\sum_{i\in I}\left( \frac{n_{t}(i)}{t}\right) M\left\Vert 
\bar{a}_{t}(i)-\bar{c}_{t}(i)\right\Vert ^{2}=M\,\mathcal{K}_{t}(\mathbf{c};%
\mathbf{f}).  \notag
\end{eqnarray}
\end{proof}

\bigskip

\begin{proof}[Proof of Theorem \protect\ref{th:proper-calibration}]
By Proposition \ref{p:K-KS}, for every proper $1$-Lipschitz scoring rule $L$
we have $\mathcal{K}_{t}^{L}(\mathbf{c})\leq \mathcal{K}_{t}(\mathbf{c}),$
and $\mathcal{K}_{t}^{L}(\mathbf{c};\Pi (\mathbf{c}))\leq \mathcal{K}_{t}(%
\mathbf{c};\Pi (\mathbf{c}))$ for every continuous binning $\Pi $.
\end{proof}

\bigskip

The existing results in the literature thus yield stochastic $\varepsilon $%
-proper-calibrated procedures (that are of type MM; and, almost
deterministic procedures of type FP), and deterministic continuously
proper-calibrated procedures (that are of type FP).

\bigskip

For instance, from Theorem 4 of Foster and Hart (2023) (with $C=\Delta $,
and thus $\gamma ^{2}=\max_{c,c^{\prime }\in \Delta }\left\Vert c-c^{\prime
}\right\Vert ^{2}=2$) we get:

\begin{theorem}
\label{th:calibration S}Let $\delta >0$ and let $\Delta _{\delta }\subset
\Delta $ be a finite $\delta $-grid of $\Delta $. Then there exists a
stochastic $\Delta _{\delta }$-forecasting procedure $\sigma $ that is $%
\delta $-proper-calibrated; specifically, 
\begin{equation*}
\mathbb{E}\left[ \mathcal{K}_{t}^{L}(\mathbf{c})\right] \leq \delta
^{2}+2|\Delta _{\delta }|\frac{\ln t+1}{t}
\end{equation*}%
for all $t\geq 1$, all sequences $\mathbf{a}_{t}\in A^{t}$, and all proper $1
$-Lipschitz scoring rules $L$ (i.e., $L\in \mathcal{L}_{1}$). Moreover, $%
\sigma $ may be taken to be $\delta $-almost deterministic (i.e., all
randomizations are $\delta $-local).
\end{theorem}

\bigskip

Theorem 11 (D) of Foster and Hart (2021) (see also Theorems 6 and 12 of
Foster and Hart 2023) yields:

\begin{theorem}
\label{th:cont}There exists a \emph{deterministic} forecasting procedure $%
\sigma $ that is continuously proper-calibrated.
\end{theorem}

\section{Proper Calibeating\label{s:proper calibeating}}

Unlike calibration, for calibeating it is \emph{not} true that
proper-calibeating is a consequence of calibeating. We show this below, and
then we provide two ways of obtaining proper calibeating: first, proving
that the simple calibeating procedure of Theorem 3 of Foster and Hart (2023)
is proper calibeating; and second, proving that calibeating the appropriate
joint binning yields proper calibeating.

\subsection{Calibeating Does Not Imply Proper Calibeating\label{s-s:example}}

We show that in general calibeating with respect to the standard quadratic
scoring rule does \emph{not} yield calibeating with respect to other proper
Lipschitz scoring rules (this is in contrast to calibration, where, as we
have seen above, calibration always entails proper calibration).

\begin{example}
\label{ex:not-S-beat}In the one-dimensional case, where $A=\{0,1\},$
consider $t=10$ periods where the actions $a_{t}$, the forecasts $b_{t}$ and 
$c_{t}$ (given in the table below as the forecasted probability of $a=1$),
are as follows:%
\begin{equation*}
\begin{tabular}{c|cccccccccc}
$t$ & $1$ & $2$ & $3$ & $4$ & $5$ & $6$ & $7$ & $8$ & $9$ & $10$ \\ 
\hline\hline
$a_{t}$ & $1$ & $0$ & $0$ & $0$ & $0$ & $1$ & $1$ & $1$ & $1$ & $0$ \\ 
$b_{t}$ & $\frac{1}{5}$ & $\frac{1}{5}$ & $\frac{1}{5}$ & $\frac{1}{5}$ & $%
\frac{1}{5}$ & $\frac{4}{5}$ & $\frac{4}{5}$ & $\frac{4}{5}$ & $\frac{4}{5}$
& $\frac{4}{5}$ \\ 
$c_{t}$ & $1$ & $0$ & $\frac{1}{2}$ & $\frac{1}{2}$ & $\frac{1}{2}$ & $\frac{%
1}{2}$ & $\frac{1}{2}$ & $\frac{1}{2}$ & $1$ & $0$%
\end{tabular}%
~~~.
\end{equation*}%
The sequences $\mathbf{b}_{t}$ and $\mathbf{c}_{t}$ are both perfectly
calibrated (i.e., $\bar{a}_{t}(b)=b$ and $\bar{a}_{t}(c)=c$ for each
forecast used), and so, for every scoring rule $L$ we have $\mathcal{K}%
_{t}^{L}(\mathbf{b})=\mathcal{K}_{t}^{L}(\mathbf{c})=0$ and\footnote{%
We use formula (\ref{eq:R-H-H}), slightly abusing notation and writing $%
H^{L}(p)$ instead of $H^{L}((p,1-p))$.} 
\begin{eqnarray*}
\mathcal{B}_{t}^{L}(\mathbf{b}) &=&\mathcal{R}_{t}^{L}(\mathbf{b})=\frac{5}{%
10}H^{L}\left( \frac{1}{5}\right) +\frac{5}{10}H^{L}\left( \frac{4}{5}%
\right) -\mathcal{H}_{t}^{L}\text{\ \ \ and} \\
\mathcal{B}_{t}^{L}(\mathbf{c}) &=&\mathcal{R}_{t}^{L}(\mathbf{c})=\frac{6}{%
10}H^{L}\left( \frac{1}{2}\right) +\frac{2}{10}H^{L}\left( 1\right) +\frac{2%
}{10}H^{L}\left( 0\right) -\mathcal{H}_{t}^{L},
\end{eqnarray*}%
where $\mathcal{H}_{t}^{L}=(5/10)H^{L}(0)+(5/10)H^{L}(1).$ For the standard
quadratic scoring rule, for which $H(p)=-(p^{2}+(1-p)^{2})$ (see Appendix %
\ref{s-a:scoring rules}) this yields%
\begin{equation}
\mathcal{B}_{t}(\mathbf{c})=\frac{3}{20}<\frac{8}{25}=\mathcal{R}_{t}(%
\mathbf{b}),  \label{eq:ex-quad}
\end{equation}%
and so $\mathbf{c}_{t}$ calibeats $\mathbf{b}_{t}$. For the $\alpha $%
-spherical scoring rule $L$ with $\alpha =2$ (which is proper and
Lipschitz), for which $H^{L}(p)=-(p^{2}+(1-p)^{2})^{1/2}$, this yields%
\footnote{%
Higher values of $\alpha $ yield higher differences between $\mathcal{B}%
_{t}^{L}(\mathbf{c})$ and $\mathcal{B}_{t}^{L}(\mathbf{b})$; for example,
for $\alpha =3$ these are $\approx 0.222$ and $\approx 0.196,$ respectively.}%
\begin{equation}
\mathcal{B}_{t}^{L}(\mathbf{c})\approx 0.1757>0.1754\approx \mathcal{R}%
_{t}^{L}(\mathbf{b}),  \label{eq:ex-L}
\end{equation}%
and so $\mathbf{c}_{t}$ does \emph{not} $L$-calibeat $\mathbf{b}_{t}$\textbf{%
.} Repeating this sequence of length $10$ periodically yields the
inequalities (\ref{eq:ex-quad}) and (\ref{eq:ex-L}) for every $t$ that is a
multiple of $10,$ and thus also in the limit as\footnote{%
Because all the scores at $t=10m+r,$ where $1\leq r\leq 9$, differ from
those at $t^{\prime }=10m$ by $O(r/t)\rightarrow 0$ as $t\rightarrow \infty $%
.} $t\rightarrow \infty $, which shows that $\mathbf{c}$ calibeats $\mathbf{b%
}$ with respect to the quadratic scoring rule but \emph{not} with respect to
the $2$-spherical scoring rule.
\end{example}

\subsection{A Simple Way to Proper Calibeat\label{s-s:simple}}

We show that the simple calibeating procedure of Theorem 3 of Foster and
Hart (2023), whereby one forecasts the current action average of the $b_{t}$%
-bin, is in fact proper-calibeating.

\begin{theorem}
\label{th:simple-calibeat-univ}Let $B$ be a finite set, and let $\zeta $ be
the deterministic $\mathbf{b}$-based forecasting procedure given by 
\begin{equation*}
c_{t}=\bar{a}_{t-1}^{\mathbf{b}}(b_{t})
\end{equation*}%
for every time $t\geq 1$ (if $t$ is the first time that $b_{t}$ is used,
take $c_{t}$ to be an arbitrary element of $\Delta $). Then the procedure $%
\zeta $ is $B$-proper-calibeating; specifically,%
\begin{equation*}
0\leq \mathcal{B}_{t}^{L}(\mathbf{c})-\mathcal{R}_{t}^{L}(\mathbf{b})\leq
2|B|\frac{\ln t+1}{t},
\end{equation*}%
for all $t\geq 1$, all sequences $\mathbf{a}_{t}\in A^{t}$ and $\mathbf{b}%
_{t}\in B^{t}$, and all proper $1$-Lipschitz scoring rules $L$ (i.e., $L\in 
\mathcal{L}_{1}$).
\end{theorem}

As in Foster and Hart (2023), we define the \emph{online }$L$-refinement
score, for a pure binning sequence $\mathbf{i}$, as follows:%
\begin{equation*}
\widetilde{\mathcal{R}}_{t}^{L}\equiv \widetilde{\mathcal{R}}_{t}^{L}(%
\mathbf{i}):=\frac{1}{t}\sum_{s=1}^{t}D(a_{s},\bar{a}_{s-1}(i_{s}))
\end{equation*}%
(take $\bar{a}_{0}(i)$ to be an arbitrary point in $\Delta $). In $\mathcal{R%
}_{t}^{L}$ one uses in each period $s$ an (offline) average of the actions, $%
\bar{a}_{t}(\cdot )$, taken over all periods from $1$ to $t$; this is
replaced in $\widetilde{\mathcal{R}}_{t}^{L}$ by the corresponding online
average of the actions, $\bar{a}_{s-1}(\cdot )$, taken only over the past
periods, from $1$ to $s-1$. We have:

\begin{proposition}
\label{p:online-R-S}Let $L$ be a proper $M$-Lipschitz scoring rule. Then%
\begin{equation*}
0\leq \widetilde{\mathcal{R}}_{t}^{L}(\mathbf{i})-\mathcal{R}_{t}^{L}(%
\mathbf{i})\leq 2M\frac{N_{t}}{t}\left( \ln \left( \frac{t}{N_{t}}\right)
+1\right) ,
\end{equation*}%
where $N_{t}:=\left\vert \{i_{s}:s\leq t\}\right\vert $ is the number of
bins used up to time $t.$
\end{proposition}

Before proving this proposition, we show how it readily proves Theorem \ref%
{th:simple-calibeat-univ}.

\begin{proof}[Proof of Theorem \protect\ref{th:simple-calibeat-univ}]
Our choice of $c_{t}=\bar{a}_{t-1}^{\mathbf{b}}(b_{t})$ makes $\mathcal{B}%
_{t}^{L}(\mathbf{c})=\widetilde{\mathcal{R}}_{t}^{L}(\mathbf{b})$ for every $%
\mathbf{a}_{t}$ and $\mathbf{b}_{t}$ and every $L$ in $\mathcal{L}$; use
Proposition \ref{p:online-R-S}, with $N_{t}\leq \left\vert B\right\vert .$
\end{proof}

\subsubsection{Online vs Offline Refinement\label{s-s-s:online refinement}}

We prove Proposition \ref{p:online-R-S}. This will follow from the following:

\begin{proposition}
\label{p:online-offline}Let $x_{1},...,x_{n}\in \Delta ,$ let $L$ be a
scoring rule, and define%
\begin{eqnarray*}
v_{n} &%
%TCIMACRO{\TeXButton{:=}{{\;:=\;}}}%
%BeginExpansion
{\;:=\;}%
%EndExpansion
&\frac{1}{n}\sum_{j=1}^{n}D(x_{j},\bar{x}_{n})\text{ and} \\
\widetilde{v}_{n} &%
%TCIMACRO{\TeXButton{:=}{{\;:=\;}}}%
%BeginExpansion
{\;:=\;}%
%EndExpansion
&\frac{1}{n}\sum_{j=1}^{n}D(x_{j},\bar{x}_{j-1}).
\end{eqnarray*}%
Then%
\begin{equation}
\widetilde{v}_{n}-v_{n}=\frac{1}{n}\sum_{j=1}^{n}jD(\bar{x}_{j},\bar{x}%
_{j-1}).  \label{eq:v-online-error}
\end{equation}%
Moreover, if $L$ is a proper $M$-Lipschitz scoring rule, then%
\begin{equation*}
0\leq \widetilde{v}_{n}-v_{n}\leq 2M\frac{\ln n+1}{n}.
\end{equation*}
\end{proposition}

\begin{proof}
Let $\xi _{n}:=n(\widetilde{v}_{n}-v_{n})$; canceling the $L(x_{j},x_{j})$
terms that appear in both sums yields 
\begin{equation*}
\xi _{n}=\sum_{j=1}^{n}L(x_{j},\bar{x}_{j-1})-\sum_{j=1}^{n}L(x_{j},\bar{x}%
_{n})=\sum_{j=1}^{n}L(x_{j},\bar{x}_{j-1})-nL(\bar{x}_{n},\bar{x}_{n})
\end{equation*}%
(for the second sum we have used $\bar{x}_{n}=(1/n)(x_{1}+...+x_{n})$). Put $%
\eta _{n}:=\xi _{n}-\xi _{n-1};$ we have%
\begin{equation*}
\eta _{n}=L(x_{n},\bar{x}_{n-1})-nL(\bar{x}_{n},\bar{x}_{n})+(n-1)L(\bar{x}%
_{n-1},\bar{x}_{n-1}).
\end{equation*}%
The sum of the first and third terms is $nL(\bar{x}_{n},\bar{x}_{n-1})$
(because $\bar{x}_{n}=(1/n)x_{n}+((n-1)/n)\bar{x}_{n-1}$), and so%
\begin{equation*}
\eta _{n}=n\left[ L(\bar{x}_{n},\bar{x}_{n-1})-L(\bar{x}_{n},\bar{x}_{n})%
\right] =nD(\bar{x}_{n},\bar{x}_{n-1}).
\end{equation*}%
Now $\xi _{n}=\sum_{j=1}^{n}\eta _{j},$ and so we have obtained the claimed
identity.

Properness gives $\eta _{j}\geq 0,$ and so $\xi _{n}\geq 0.$ The Lipschitz
condition gives, by Lemma \ref{l:O(diff^2)},%
\begin{equation*}
\eta _{j}\leq jM\left\Vert \bar{x}_{j}-\bar{x}_{j-1}\right\Vert
^{2}=jM\left\Vert \frac{1}{j}(\bar{x}_{j-1}-x_{j})\right\Vert ^{2}\leq jM%
\frac{2}{j^{2}}=\frac{2M}{j}
\end{equation*}%
(we used $\left\Vert x-y\right\Vert ^{2}\leq 2$ for all $x,y\in \Delta $).
Therefore%
\begin{equation*}
\xi _{n}=\sum_{j=1}^{n}\eta _{j}\leq 2M\sum_{j=1}^{n}\frac{1}{j}\leq 2M(\ln
n+1),
\end{equation*}%
completing the proof.
\end{proof}

\bigskip

\begin{proof}[Proof of Proposition \protect\ref{p:online-R-S}]
For each bin $i$ with $n_{t}(i)>0,$ Proposition \ref{p:online-offline} gives%
\begin{equation*}
0\leq \frac{1}{n_{t}(i)}\sum_{s\leq t:i_{s}=i}D(a_{s},\bar{a}_{s-1}(i))-%
\frac{1}{n_{t}(i)}\sum_{s\leq t:i_{s}=i}D(a_{s},\bar{a}_{t}(i))\leq 2M\frac{%
\ln n_{t}(i)+1}{n_{t}(i)}.
\end{equation*}%
Averaging over all these $i$ with weights $n_{t}(i)/t$ then yields%
\begin{equation*}
0\leq \widetilde{\mathcal{R}}_{t}^{L}-\mathcal{R}_{t}^{L}\leq 2M\frac{1}{t}%
\sum_{i}(\ln n_{t}(i)+1).
\end{equation*}%
The maximum of the righthand side is attained when all the positive $%
n_{t}(i) $ are equal (because $\ln $ is concave), i.e., when $%
n_{t}(i)=t/N_{t}$ (because there are $N_{t}$ nonempty bins, and the sum of
all the $n_{t}(i)$ is $t$); this yields the claimed bound.
\end{proof}

\bigskip

\noindent \textbf{Remarks.} \emph{(a)} In the quadratic case formula (\ref%
{eq:v-online-error}) yields Proposition 2 of Foster and Hart 2023 on the
online variance.

\emph{(b)} In the logarithmic case one may use the regularization of adding
a positive constant to each bin (see Appendix A.9 in Foster and Hart 2022).
However, this does \emph{not} work when the slope at the boundary of $\Delta 
$ is much higher. For instance, take the $\alpha $-power scoring rule with $%
\alpha =-1,$ binary ($m=2$), and consider a single bin. Let the sequence of
actions be $x_{1}=(1,0)$ and then $x_{j}=(0,1)$ forever (i.e., for all $%
j\geq 2$), then%
\begin{equation*}
\eta _{j}=jD(\bar{x}_{j},\bar{x}_{j-1})=1+\frac{1}{(j-1)(j-2)^{2}}\geq 1
\end{equation*}%
for all $j\geq 3,$ which shows that\footnote{%
Ignore the first $2$ periods where $\delta _{i}$ is infinite---these are the
\textquotedblleft regularization" periods whereby we \textquotedblleft seed"
the bin---and start counting only from $i=3.$} $\widetilde{v}_{n}-v_{n}\geq
(n-2)/n\rightarrow 1,$ and the online refinement is at a distance of about $1
$ from the offline refinement.

\subsection{Proper Calibeating by a Proper Calibrated Procedure\label%
{s-s:joint}}

Here we prove that calibeating the appropriate joint binning, as in Theorem
5 of Foster and Hart (2023), yields proper-calibeating, moreover, by a
proper-calibrated procedure.

\begin{theorem}
\label{th:beat-by-calib S}Let $B$ be a finite set, and let $\Delta _{\delta
}\subset \Delta $ be a finite $\delta $-grid of $\Delta $ for some $\delta >0
$. Then there exists a stochastic $\mathbf{b}$-based $\Delta _{\delta }$%
-forecasting procedure $\zeta $ that is $(\delta ,B)$-proper-calibeating and 
$\delta $-proper-calibrated; specifically,%
\begin{equation*}
\mathbb{E}\left[ \mathcal{B}_{t}^{L}(\mathbf{c})-\mathcal{R}_{t}^{L}(\mathbf{%
b}\times \mathbf{c})\right] \leq \delta ^{2}+2|B|\,|\Delta _{\delta }|\frac{%
\ln t+1}{t},
\end{equation*}%
and thus,%
\begin{eqnarray*}
\mathbb{E}\left[ \mathcal{B}_{t}^{L}(\mathbf{c})-\mathcal{R}_{t}^{L}(\mathbf{%
b})\right]  &\leq &\delta ^{2}+2|B|\,|\Delta _{\delta }|\frac{\ln t+1}{t}%
\text{\ \ and} \\
\mathbb{E}\left[ \mathcal{K}_{t}^{L}(\mathbf{c})\right] =\mathbb{E}\left[ 
\mathcal{B}_{t}^{L}(\mathbf{c})-\mathcal{R}_{t}^{L}(\mathbf{c})\right] 
&\leq &\delta ^{2}+2|B|\,|\Delta _{\delta }|\frac{\ln t+1}{t}
\end{eqnarray*}%
for all $t\geq 1$ and all sequences $\mathbf{a}_{t}\in A^{t}$ and $\mathbf{b}%
_{t}\in B^{t}$, and all proper $1$-Lipschitz scoring rules $L$ (i.e., $L\in 
\mathcal{L}_{1}$). Moreover, $\zeta $ may be taken to be $\delta $-almost
deterministic.
\end{theorem}

\begin{proof}
The decomposition Theorem \ref{th:S-decompose} (applied to both $\mathcal{B}%
_{t}^{L}$ and $\mathcal{B}_{t}$) together with Proposition \ref{p:K-KS} yield%
\begin{eqnarray*}
\mathcal{B}_{t}^{L}(\mathbf{c})-\mathcal{R}_{t}^{L}(\mathbf{b}\times \mathbf{%
c}) &=&\mathcal{K}_{t}^{L}(\mathbf{c};\mathbf{b}\times \mathbf{c}) \\
&\leq &\mathcal{K}_{t}(\mathbf{c};\mathbf{b}\times \mathbf{c})=\mathcal{B}%
_{t}(\mathbf{c})-\mathcal{R}_{t}(\mathbf{b}\times \mathbf{c}).
\end{eqnarray*}%
Apply Theorem 5 of Foster and Hart (2023) with $C=\Delta $ (and thus $\gamma
^{2}=\max_{c,d\in \Delta }\left\Vert c-d\right\Vert ^{2}=2$) to get the
first inequality, and then Proposition \ref{p:refine-S} below, which shows
that coarsening the binning can only increase the refinement score, to get
the other two inequalities.
\end{proof}

\bigskip \bigskip 

\noindent \textbf{Remark. }The following four statements regarding $\mathbf{c%
}$ and the joint binning $\mathbf{b}\times \mathbf{c}$ are equivalent:

\begin{description}
\item[(J1)] $\mathbf{c}$ calibeats $\mathbf{b\times c}$, i.e., $\mathcal{B}(%
\mathbf{c})\leq \mathcal{R}(\mathbf{b}\times \mathbf{c})$; equivalently,%
\footnote{%
Because we always have $\mathcal{R}(\mathbf{c})\leq \mathcal{B}(\mathbf{c})$%
, and $\mathcal{R}(\mathbf{b}\times \mathbf{c})\leq \mathcal{R}(\mathbf{c})$
(since the $\mathbf{b}\times \mathbf{c}$-binning is a refinement of the $%
\mathbf{c}$-binning). Similarly for every scoring rule $L$, i.e., for (J4)
below.} $\mathcal{B}(\mathbf{c})=\mathcal{R}(\mathbf{c})=\mathcal{R}(\mathbf{%
b}\times \mathbf{c})$.

\item[(J2)] $\mathbf{c}$ is calibrated on the $\mathbf{b\times c}$-binning,
i.e., $\mathcal{K}(\mathbf{c};\mathbf{b}\times \mathbf{c})=0$.

\item[(J3)] $\mathbf{c}$ is proper-calibrated on the $\mathbf{b\times c}$%
-binning, i.e., $\mathcal{K}^{L}(\mathbf{c};\mathbf{b}\times \mathbf{c})=0$
for every $L$ in $\mathcal{L}$.

\item[(J4)] $\mathbf{c}$ proper-calibeats $\mathbf{b\times c}$, i.e., $%
\mathcal{B}^{L}(\mathbf{c})\leq \mathcal{R}^{L}(\mathbf{b}\times \mathbf{c})$
for every $L$ in $\mathcal{L}$; equivalently, $\mathcal{B}^{L}(\mathbf{c})=%
\mathcal{R}^{L}(\mathbf{c})=\mathcal{R}^{L}(\mathbf{b}\times \mathbf{c})$
for every $L$ in $\mathcal{L}$.
\end{description}

\noindent Indeed, since the $\mathbf{b}\times \mathbf{c}$-binning is a
refinement of the $\mathbf{c}$-binning, we get the decomposition $\mathcal{B}%
^{L}(\mathbf{c})=\mathcal{R}^{L}(\mathbf{b}\times \mathbf{c})+\mathcal{K}%
^{L}(\mathbf{c};\mathbf{b}\times \mathbf{c})$, which immediately yields (J1) 
$\Longleftrightarrow $ (J2), and (J3) $\Longleftrightarrow $ (J4). As for
(J2) $\Longleftrightarrow $ (J3), it follows from $\mathcal{K}^{L}\leq M\,%
\mathcal{K}$ for every $L$ in $\mathcal{L}$, and the standard quadratic
scoring rule being in $\mathcal{L}$. Moreover, (J1) implies that $\mathbf{c}$
is calibrated (because $\mathcal{K}(\mathbf{c})=\mathcal{B}(\mathbf{c})-%
\mathcal{R}(\mathbf{c})=0$), and (J4) that it is proper-calibrated (because $%
\mathcal{K}^{L}(\mathbf{c})=\mathcal{B}^{L}(\mathbf{c})-\mathcal{R}^{L}(%
\mathbf{c})=0$).

\bigskip 

\noindent \textbf{Remark. }Calibeating together with calibration does \emph{%
not} suffice to get $L$-calibeating; one needs $\mathbf{c}$ to \emph{%
calibeat the joint} $\mathbf{b}\times \mathbf{c}$ of $\mathbf{b}$ and $%
\mathbf{c}$. See Example \ref{ex:not-S-beat}, where $L$ is the proper
Lipschitz $2$-spherical scoring rule: $\mathbf{c}$ is calibrated (and thus $L
$-calibrated) and calibeats $\mathbf{b}$, but it does not $L$-calibeat $%
\mathbf{b}$. Indeed, $\mathbf{c}$ is not calibrated with respect to the $%
\mathbf{b}\times \mathbf{c}$-binning: the $(b=1/5,c=1/2)$-bin is not $%
\mathbf{c}$-calibrated, as the average action there is $0$ and not $1/2$. As
seen in the above proof, to get $L$-calibeating from calibeating we rely on
the decomposition Theorem \ref{th:S-decompose} which requires the binning to
be a refinement of the $\mathbf{c}$-binning, and so it does not apply to an
arbitrary $\mathbf{b}$-binning,\footnote{%
For our simple calibeating procedure we use a different tool, namely,
Proposition \ref{p:online-R-S}.} but it does apply to the joint $\mathbf{b}$%
\textbf{$\times $}$\mathbf{c}$-binning. Appendix \ref{s-a:bxc} provides a
setup where we vary only the frequencies of the bins, suggesting that a
\textquotedblleft natural\textquotedblright\ proof of proper-calibeating
from calibeating together with calibration might well require calibeating
the joint.

\subsubsection{Refined Refinement\label{s-s-s:refined refinement}}

For the proof we need to generalize to proper scoring rules the result that
the refinement score can only decrease when the binning becomes finer (for
the quadratic scoring rule, this is Proposition 11 of Foster and Hart 2023).
While in this section we need it only for pure binnings, we state it already
for the more general fractional binnings (to be used in the following
section).

Let $\mathbf{f}$ be a fractional binning on a set of bins $I,$ and $\mathbf{g%
}$ a fractional binning sequence on a set of bins $J$ (thus $f_{s}\in \Delta
(I)$ and $g_{s}\in \Delta (J)$). We say that $\mathbf{f}$ is a \emph{%
refinement} of $\mathbf{g}$ (or $\mathbf{g}$ is a \emph{coarsening} of $%
\mathbf{f}$) if each $j$\textbf{-}bin is a union of $i$-bins, with $g_{s}(j)$
the sum of the corresponding $f_{s}(i)$; i.e., there is a partition $I=\cup
_{j\in J}I(j)$ of $I$ into disjoint sets $I(j)$ for $j\in J,$ and $%
g_{s}(j)=\sum_{i\in I(j)}f_{s}(i)$ for every $j\in J$ and $s\geq 1$.

\begin{proposition}
\label{p:refine-S}If the binning sequence $\mathbf{f}$ is a refinement of
the binning sequence $\mathbf{g}$, then%
\begin{equation*}
\mathcal{R}_{t}^{L}(\mathbf{f})\leq \mathcal{R}_{t}^{L}(\mathbf{g})
\end{equation*}%
for every proper scoring rule $L$ and every $t\geq 1.$
\end{proposition}

\begin{proof}
Let $I$ and $J$ be the sets of bins of $\mathbf{f}$ and $\mathbf{g}$,
respectively. It suffices to prove the claim when $J$ has only one bin; we
then apply it to each $j$-bin separately and average over $j$ to get the
general result. Let $\bar{a}_{t}$ denote the overall average of the actions
(i.e., the average in the single bin in $J$), then we have%
\begin{eqnarray*}
\frac{1}{t}\sum_{i\in I}\sum_{s=1}^{t}f_{s}(i)L(a_{s},\bar{a}_{t}(i))
&=&\sum_{i\in I}\left( \frac{n_{t}(i)}{t}\right) \sum_{s=1}^{t}\left( \frac{%
f_{s}(i)}{n_{t}(i)}\right) L(a_{s},\bar{a}_{t}(i)) \\
&=&\sum_{i\in I}\left( \frac{n_{t}(i)}{t}\right) L(\bar{a}_{t}(i),\bar{a}%
_{t}(i))=\sum_{i\in I}\left( \frac{n_{t}(i)}{t}\right) H(\bar{a}_{t}(i)) \\
&\leq &H\left( \sum_{i}\frac{n_{t}(i)}{t}\bar{a}_{t}(i)\right) =H(\bar{a}%
_{t}) \\
&=&L(\bar{a}_{t},\bar{a}_{t})=\frac{1}{t}\sum_{s=1}^{t}L(a_{s},\bar{a}_{t}),
\end{eqnarray*}%
where the inequality is by the concavity of the function $H$ (and we have
for convenience dropped the subscript $t$ from $\bar{a}_{t}(i)$ and $%
n_{t}(i) $). Subtracting $\mathcal{H}_{t}^{L}$ from both sides yields the
desired inequality.
\end{proof}

\bigskip

\begin{corollary}
\label{c:refine-S}If the binning sequence $\mathbf{f}$ is a refinement of
the binning sequence $\mathbf{g}$ and the binning sequence $\mathbf{g}$ is a
refinement of the binning sequence $\mathbf{c}$ then%
\begin{equation*}
\mathcal{K}_{t}^{L}(\mathbf{c};\mathbf{f})\geq \mathcal{K}_{t}^{L}(\mathbf{c}%
;\mathbf{g})\geq \mathcal{K}_{t}^{L}(\mathbf{c})
\end{equation*}%
for every proper scoring rule $L$ and every $t\geq 1.$
\end{corollary}

\begin{proof}
Since both $\mathbf{f}$ and $\mathbf{g}$ refine $\mathbf{c}$, the
decomposition of \ref{th:S-decompose} applies to each one of $\mathbf{f}$, $%
\mathbf{g}$, and $\mathbf{c}$, and so $\mathcal{R}^{L}(\mathbf{f})+\mathcal{K%
}^{L}(\mathbf{c};\mathbf{f})\mathbf{=}\mathcal{R}^{L}(\mathbf{g})+\mathcal{K}%
^{L}(\mathbf{c};\mathbf{g})=\mathcal{R}^{L}(\mathbf{c})+\mathcal{K}_{t}^{L}(%
\mathbf{c};\mathbf{c})$ (they all equal $\mathcal{B}^{L}(\mathbf{c})$);
apply Proposition \ref{p:refine-S}.
\end{proof}

\subsection{Proper Calibeating by a Deterministic Continuously Proper
Calibrated Procedure\label{s-s:continuous}}

The result of this section is the proper counterpart of Theorems 6 and 12 of
Foster and Hart (2023, 2026). Since fractional binnings in general do not
refine the standard by-forecast binning (because each bin may well contain
forecasts with different values), and the decomposition $\mathcal{B}^{L}=%
\mathcal{K}^{L}+\mathcal{R}^{L}$ is no longer valid, we use Theorem \ref%
{th:decompose-delta} instead.

\begin{theorem}
\label{th:calibeat by cont cal S}Let $B$ be a finite set. Then there exists
a \emph{deterministic} $\mathbf{b}$-based forecasting procedure $\zeta $
that is $B$-proper-calibeating, and continuously proper-calibrated.
Specifically: first, for every continuous binning $\Pi $ there is a
deterministic $\mathbf{b}$-based forecasting procedure $\zeta $such that%
\begin{equation}
\mathcal{K}_{t}^{L}(\mathbf{c};\mathbf{b}\times \Pi (\mathbf{c}))\leq o(1);
\label{eq:34}
\end{equation}%
and second, for the continuous binning $\Pi ^{\ast }$ of Foster and Hart
(2026), (\ref{eq:34}) implies that 
\begin{equation*}
\mathcal{B}_{t}^{L}(\mathbf{c})\leq \mathcal{R}_{t}^{L}(\mathbf{b})+o(1),
\end{equation*}%
and that $\zeta $ is continuously $L$-calibrated. All these hold as $%
t\rightarrow \infty $ uniformly over all sequences $\mathbf{a}$ and $\mathbf{%
b}$ and scoring rules $L$ in $\mathcal{L}_{1}$.
\end{theorem}

\begin{proof}
For every continuous binning $\Pi $ the procedure of Theorem 12 in Foster
and Hart (2026) yields (\ref{eq:34}) for the quadratic scoring rule, and
thus uniformly for all $L$ in $\mathcal{L}$ by Proposition \ref{p:K-KS}.

Consider now $\Pi ^{\ast }$, which contains a sequence $(\Pi _{n})_{n\geq 1}$
of $\delta _{n}$-local continuous binnings with $\delta _{n}>0$ converging
to $0.$ As shown in the proof of Theorem 12 in Foster and Hart (2026), (\ref%
{eq:34}) for $\Pi ^{\ast }$ implies that for each $n\geq 1$ we have%
\begin{equation*}
\mathcal{K}_{t}(\mathbf{c};\mathbf{b}\times \Pi _{n}(\mathbf{c}))\leq o(1)
\end{equation*}%
(see (E1)\ there). Using the approximate decomposition of Theorem \ref%
{th:decompose-delta} (the binning $\Pi _{n}(\mathbf{c})$ is $\delta _{n}$%
-local with respect to $\mathbf{c}$, and thus so is its refinement $\mathbf{b%
}\times \Pi _{n}(\mathbf{c})$) and, again, Proposition \ref{p:K-KS}, we get%
\begin{eqnarray*}
\mathcal{B}_{t}^{L}(\mathbf{c})-\mathcal{R}_{t}^{L}(\mathbf{b}\times \Pi
_{n}(\mathbf{c})) &\leq &\mathcal{K}_{t}^{L}(\mathbf{c};\mathbf{b}\times \Pi
_{n}(\mathbf{c}))+2\delta _{n} \\
&\leq &\mathcal{K}_{t}(\mathbf{c};\mathbf{b}\times \Pi _{n}(\mathbf{c}%
))+2\delta _{n}\leq 2\delta _{n}+o(1)
\end{eqnarray*}%
(uniformly for all $L$ in $\mathcal{L}$). Since $\mathcal{R}_{t}^{L}(\mathbf{%
b}\times \Pi _{n}(\mathbf{c}))\leq \mathcal{R}_{t}^{L}(\mathbf{b})$ by
Proposition \ref{p:refine-R}, we get%
\begin{equation*}
\mathcal{B}_{t}^{L}(\mathbf{c})-\mathcal{R}_{t}^{L}(\mathbf{b})\leq 2\delta
_{n}+o(1).
\end{equation*}%
Therefore $\mathcal{B}_{t}^{L}(\mathbf{c})-\mathcal{R}_{t}^{L}(\mathbf{b}%
)\leq 3\delta _{n}$ for all $t$ large enough; since $\delta _{n}\rightarrow
0,$ this yields $\mathcal{B}_{t}^{L}(\mathbf{c})-\mathcal{R}_{t}^{L}(\mathbf{%
b})\leq o(1).$
\end{proof}

\subsection{Proper Multicalibeating\label{s-s:multi}}

Suppose that there are $N\geq 1$ forecasting sequences, $\mathbf{b}%
^{n}=(b_{t}^{n})_{t\geq 1}$ for $n=1,2,...,N$. We assume that each $\mathbf{b%
}^{n}$ uses only finitely many forecasts: there is a finite set $B^{n}$ such
that $b_{t}^{n}\in B^{n}$ for all $t\geq 1$. Put $\mathbf{b=}(\mathbf{b}%
^{1},...,\mathbf{b}^{N});$ we are looking for a $\mathbf{b}$-based
forecasting procedure---i.e., $c_{t}$ is determined after all the $%
b_{t}^{1},...,b_{t}^{N}$ are announced (and hence is a function of $\mathbf{a%
}_{t-1},\mathbf{c}_{t-1},\mathbf{b}_{t}^{1},...,\mathbf{b}_{t}^{N}$---that
simultaneously calibeats all the $\mathbf{b}^{n}$ sequences. By applying the
results of the previous section to the joint binning $\mathbf{b}_{1}\times 
\mathbf{b}_{2}\times ...\times \mathbf{b}_{N}$ we get the counterpart of
Theorem 7 in Foster and Hart (2023):

\begin{theorem}
\label{th:multi}Let $B^{1},B^{2},...,B^{N}$ be finite sets. Then:

\begin{description}
\item[(i)] There exists a simple deterministic $(\mathbf{b}^{1},...,\mathbf{b%
}^{N})$-based forecasting procedure $\zeta $ that is $B^{n}$%
-proper-calibeating for all $n=1,...,N;$ specifically, the forecast of $%
\zeta $ in period $t$ is $c_{t}=\bar{a}_{t-1}^{\mathbf{b}^{1},...,\mathbf{b}%
^{N}}(b_{t}^{1},...,b_{t}^{N})$, the average of the actions in all past
periods $s\leq t-1$ where the combination $(b_{t}^{1},...,b_{t}^{N})$ was
used (if $t$ is the first period in which $(b_{t}^{1},...,b_{t}^{N})$ is
used, take $c_{t}\in \Delta $ to be arbitrary).

\item[(ii)] For every finite $\delta $-grid $\Delta _{\delta }$ of $\Delta $
there exists a stochastic $(\mathbf{b}^{1},...,\mathbf{b}^{N})$-based $%
\Delta _{\delta }$-forecasting procedure $\zeta $ that is $(\delta ,B^{n})$%
-proper-calibeating for all $n=1,...,N$ and is $\delta $-proper-calibrated.
Moreover, $\zeta $ may be taken to be $\delta $-almost deterministic.

\item[(iii)] There exists a deterministic $(\mathbf{b}^{1},...,\mathbf{b}%
^{N})$-based $\Delta $-forecasting procedure $\zeta $ that is $B^{n}$%
-proper-calibeating for all $n=1,...,N$ and is continuously
proper-calibrated.
\end{description}
\end{theorem}

\bigskip

\appendix{}

\section{Scoring Rules\label{s-a:scoring rules}}

A \emph{scoring rule} $L_{A}:A\times \Delta \rightarrow \mathbb{R}$
associates to every forecast $c\in \Delta $ and every realized action $a\in A
$ a loss\footnote{%
There are scoring rules, such as the logarithmic scoring rule, that allow
the loss to be infinite; as they are not Lipschitz-continuous, we will not
deal with them here (for \textquotedblleft
log-calibeating,\textquotedblright\ see Appendix A.9 in [FH2022]).} $%
L_{A}(a,c)$. The function $L_{A}$ is linearly extended to $L:\Delta \times
\Delta \rightarrow \mathbb{R}$ by $L(d,c):=\sum_{a\in A}d(a)L_{A}(a,c)$;
thus, $L(d,c)=\mathbb{E}_{a\sim d}\left[ L_{A}(a,c)\right] $ is the expected
loss when the realization is distributed according to\footnote{%
For $a$ in A we thus have $L_{A}(a,c)=L(\mathbf{1}_{a},c).$} $d\in \Delta .$
Writing $\mathbf{L}(c)$ for the vector $(L_{A}(a,c))_{a\in A}$ in $\mathbb{R}%
^{A}$, we thus have%
\begin{equation*}
L(d,c)=d\cdot \mathbf{L}(c)
\end{equation*}%
for every $c,d\in \Delta .$

A scoring rule $L$ is \emph{proper} if $L(d,c)\geq L(d,d)$ for every $c,d\in
\Delta $; i.e., forecasting the correct distribution minimizes the expected
loss. If $L(d,c)>L(d,d)$ for all $c\neq d$ then $L$ is \emph{strictly proper.%
}

The following is well known (see ...).

\begin{proposition}
\label{p:proper}A scoring rule $L$ is proper if and only if there exists a
concave function $H:\Delta \rightarrow \mathbb{R}$ and a supergradient
selection $\mathbf{G}:\Delta \rightarrow \mathbb{R}^{A}$ (i.e., for every $%
c\in \Delta $ the vector $\mathbf{G}(c)$ is a supergradient of $H$ at%
\footnote{%
I.e., $H(d)\leq H(c)+(d-c)\cdot \mathbf{G}(c)$ for every $d\in \Delta .$} $c$%
) such that%
\begin{equation}
L(d,c)=H(c)+(d-c)\cdot \mathbf{G}(c)  \label{eq:S from H}
\end{equation}%
for every $c,d\in \Delta .$
\end{proposition}

\begin{proof}
Assume that $L$ is proper. Then $H(d):=L(d,d)=\min_{c\in \Delta
}L(d,c)=\min_{c\in \Delta }d\cdot \mathbf{L}(c)$ is the minimum of linear
functions of $d$, and thus $H$ is concave. The vector $\mathbf{L}(c)$ is a
supergradient of $H$ at $c,$ because $H(d)-H(c)=L(d,d)-L(c,c)\leq
L(d,c)-L(c,c)=d\cdot \mathbf{L}(c)-c\cdot \mathbf{L}(c)=(d-c)\cdot \mathbf{L}%
(c),$ where the inequality is by properness.

Conversely, given a concave $H$ with a supergradient selection $\mathbf{G}$,
the supergradient inequality $H(d)\leq H(c)+(d-c)\cdot \mathbf{G}(c)$
becomes $L(d,d)\leq L(d,c)$ for the function $L$ that is defined by (\ref%
{eq:S from H}), and so $L$ is proper.
\end{proof}

\bigskip

The concave function $H\equiv H^{L}$ given by $H(c):=L(c,c)$ for a proper
scoring rule $L$ is usually referred to as the $L$-\emph{entropy}.

\bigskip

\noindent \textbf{Remark. }To resolve some confusion here: in the second
part of the proof the supergradient $\mathbf{G}(c)$ need \emph{not} be the
vector $\mathbf{L}(c)$ that was used in the first part. The reason is that
if $\mathbf{G}(c)$ is a supergradient then so is $\mathbf{G}(c)+\lambda e$
for any real $\lambda $ (where $e=(1,...,1)\in \mathbb{R}^{A}$), because on
the domain $\Delta $ we have $(d-c)\cdot \lambda e=0$ for every $c,d\in
\Delta $ (and so formula (\ref{eq:S from H}) is not affected by adding $%
\lambda e$ to $\mathbf{G}(c)$).\footnote{%
Even if $H$ is differentiable on the entire space $\mathbb{R}^{A}$, the
vector $\mathbf{G}(c)$ may well be different from the gradient $\nabla H(c)$%
; however, we always have $(d-c)\cdot \mathbf{G}(c)=(d-c)\cdot \nabla H(c)$
for $c,d$ in $\Delta .$}

\bigskip

Define the $L$\emph{-divergence}\footnote{%
At times this is written $D(d||c)$ (as is the case for the Kullback-Leibler
divergence).} $D\equiv D^{L}:\Delta \times \Delta \rightarrow \mathbb{R}$ by%
\begin{equation*}
D(d,c)\,:=\,L(d,c)-L(d,d).
\end{equation*}%
Then $L$ is proper if and only if $D$ is always $\geq 0.$ Geometrically, in
this case we have $D(d,c)=H(c)+(d-c)\cdot \mathbf{G}(c)-H(d)$ by (\ref{eq:S
from H}), and so the divergence $D(d,c)$ is equal to how much the tangent to
the concave function $H$ at $c$ (with slope $\mathbf{G}(c)$) is above $H$ at
the point $d$ (the \textquotedblleft Bregman divergence\textquotedblright ).

A scoring rule $L$ is \emph{Lipschitz }continuous if the functions $%
L_{A}(a,\cdot )$ for all $a\in A,$ and thus $L(d,\cdot )$ for all $d\in
\Delta ,$ are Lipschitz;\footnote{%
This obtains when the entropy $H$ is a so-called \textquotedblleft
L-smooth\textquotedblright\ function, which means that it is differentiable
and its gradient is Lipschitz.} it is convenient to require this directly on
the vector function $\mathbf{L}$: we say that $L$ is $M$\emph{-Lipschitz}
for some finite $M$ if%
\begin{equation*}
\left\Vert \mathbf{L}(c)-\mathbf{L}(c^{\prime })\right\Vert \leq M\left\Vert
c-c^{\prime }\right\Vert
\end{equation*}%
for all $c,c^{\prime }\in \Delta $, where $\left\Vert \cdot \right\Vert $
denotes the standard Euclidean norm.\footnote{%
Requiring instead that each $L_{A}(a,\cdot )$ be $M$-Lipschitz would add an
unwieldy factor of $\sqrt{\left\vert A\right\vert }$ throughout.} We record
two immediate but useful implications of the Lipschitz condition.

\begin{lemma}
\label{l:O(diff^2)}Let $L$ be an $M$-Lipschitz scoring rule. Then

\begin{description}
\item[(i)] For every $c,c^{\prime },d$ in $\Delta $ we have%
\begin{equation*}
\left\vert D(d,c)-D(d,c^{\prime })\right\vert =\left\vert
L(d,c)-L(d,c^{\prime })\right\vert \leq M\left\Vert c-c^{\prime }\right\Vert
.
\end{equation*}

\item[(ii)] For every $c,d$ in $\Delta $ we have%
\begin{equation*}
D(d,c)+D(c,d)\leq M\left\Vert c-d\right\Vert ^{2},
\end{equation*}%
and so, if $L$ is proper then%
\begin{equation*}
0\leq D(d,c)\leq M\left\Vert c-d\right\Vert ^{2}.
\end{equation*}
\end{description}
\end{lemma}

\begin{proof}
Using (\ref{eq:d*Lambda(c)}), the definition of $D$, and $\left\Vert
d\right\Vert \leq 1$ for $d\in D$, gives:%
\begin{eqnarray*}
\left\vert D(d,c)-D(d,c^{\prime })\right\vert &=&\left\vert
L(d,c)-L(d,c^{\prime })\right\vert =\left\vert d\cdot (\mathbf{L}(c)-\mathbf{%
L}(c^{\prime }))\right\vert \\
&\leq &\left\Vert d\right\Vert \left\Vert \mathbf{L}(c)-\mathbf{L}(c^{\prime
})\right\Vert \leq M\left\Vert c-c^{\prime }\right\Vert
\end{eqnarray*}

and%
\begin{eqnarray*}
D(d,c)+D(c,d) &=&d\cdot \mathbf{L}(c)-d\cdot \mathbf{L}(d)+c\cdot \mathbf{L}%
(d)-c\cdot \mathbf{L}(c) \\
&=&(c-d)\cdot (\mathbf{L}(d)-\mathbf{L}(c)) \\
&\leq &\left\Vert c-d\right\Vert \left\Vert \mathbf{L}(d)-\mathbf{L}%
(c)\right\Vert \leq M\left\Vert c-d\right\Vert ^{2}.
\end{eqnarray*}
\end{proof}

\subsection{Examples of Proper Lipschitz Scoring Rules\label{s-s-a:scoring
examples}}

We provide a number of classical examples of proper scoring rules; they are
all Lipschitz, and strictly proper.

\bigskip

\begin{itemize}
\item \emph{Quadratic} (Brier): 
\begin{eqnarray*}
L_{A}(a,c) &=&-2c(a)+\left\Vert c\right\Vert ^{2}, \\
L(d,c) &=&-2c\cdot d+\left\Vert c\right\Vert ^{2}, \\
H(c) &=&-\left\Vert c\right\Vert ^{2}, \\
D(d,c) &=&\left\Vert c-d\right\Vert ^{2}.
\end{eqnarray*}

\item $\alpha $-\emph{Spherical} for $\alpha \geq 2$: 
\begin{eqnarray*}
L_{A}(a,c) &=&-\frac{c(a)^{\alpha -1}}{\left( \sum_{a}c(a)^{\alpha }\right)
^{\frac{\alpha -1}{\alpha }}}, \\
L(d,c) &=&-\frac{\sum_{a}d(a)c(a)^{\alpha -1}}{\left( \sum_{a}c(a)^{\alpha
}\right) ^{\frac{\alpha -1}{\alpha }}}, \\
H(c) &=&-\left( \sum_{a}c(a)^{\alpha }\right) ^{\frac{1}{\alpha }%
}=-\left\Vert c\right\Vert _{\alpha } \\
D(d,c) &=&-\frac{\sum_{a}d(a)c(a)^{\alpha -1}}{\left( \sum_{a}c(a)^{\alpha
}\right) ^{\frac{\alpha -1}{\alpha }}}+\left( \sum_{a}d(a)^{\alpha }\right)
^{\frac{1}{\alpha }}.
\end{eqnarray*}%
(for $1<\alpha <2$ the scoring rule is proper but not Lipschitz).

\item $\alpha $-\emph{Power} (\emph{Tsallis}) for $\alpha \geq 2$:%
\begin{eqnarray*}
L_{A}(a,c) &=&-\frac{1}{\alpha -1}c(a)^{\alpha -1}+\frac{1}{\alpha }%
\sum_{a}c(a)^{\alpha }, \\
L(d,c) &=&-\frac{1}{\alpha -1}\sum_{a}d(a)c(a)^{\alpha -1}+\frac{1}{\alpha }%
\sum_{a}c(a)^{\alpha }, \\
H(c) &=&-\frac{1}{\alpha (\alpha -1)}\sum_{a}c(a)^{\alpha }=-\frac{1}{\alpha
(\alpha -1)}\left( \left\Vert c\right\Vert _{\alpha }\right) ^{\alpha } \\
D(d,c) &=&\frac{1}{\alpha }\sum_{a}c(a)^{\alpha }+\frac{1}{\alpha (\alpha -1)%
}\sum_{a}d(a)^{\alpha }-\frac{1}{\alpha -1}\sum_{a}d(a)c(a)^{\alpha -1}
\end{eqnarray*}%
(for $\alpha =2$ this gives the quadratic score divided by $2$; the rule is
proper for all $\alpha \neq 0,1$, and Lipschitz only for $a\geq 2$).
\end{itemize}

\section{Towards a Natural Proof of Proper Calibeating by a Calibrated
Procedure\label{s-a:bxc}}

We show here that our line of proof of proper-calibeating by a
proper-calibrated procedure, which assumes calibeating the joint binning,
may well be necessary. This is suggested by keeping all the bin averages
fixed, and allowing the relative frequencies of the various bins to vary. We
strat with a general result, and then apply it to our setup.

\subsection{A General Result\label{s-s-a:bxc general}}

Let $I$ and $J$ be finite sets, $X=(x_{ij})_{i\in I,j\in J}$ a matrix whose
entries are $m$-dimensional real vectors, i.e., $x_{ij}\in \mathbb{R}^{m}$,
and let $W=(w_{ij})_{i\in I,j\in J}$ be a weight matrix, i.e., $w_{ij}\geq 0$
for all $i,j,$ and $\sum_{i\in I}\sum_{j\in J}w_{ij}=1$. Define%
\begin{eqnarray*}
w_{i\cdot } &%
%TCIMACRO{\TeXButton{:=}{{\;:=\;}}}%
%BeginExpansion
{\;:=\;}%
%EndExpansion
&\sum_{j\in J}w_{ij} \\
w_{\cdot j} &%
%TCIMACRO{\TeXButton{:=}{{\;:=\;}}}%
%BeginExpansion
{\;:=\;}%
%EndExpansion
&\sum_{i\in I}w_{ij} \\
r_{i} &%
%TCIMACRO{\TeXButton{:=}{{\;:=\;}}}%
%BeginExpansion
{\;:=\;}%
%EndExpansion
&\bar{x}_{i\cdot }=\sum_{j\in J}\frac{w_{ij}}{w_{i\cdot }}x_{ij} \\
c_{j} &%
%TCIMACRO{\TeXButton{:=}{{\;:=\;}}}%
%BeginExpansion
{\;:=\;}%
%EndExpansion
&\bar{x}_{\cdot j}=\sum_{i\in I}\frac{w_{ij}}{w_{\cdot j}}x_{ij}
\end{eqnarray*}%
(thus, $w_{i\cdot }$ and $w_{\cdot j}$ are the marginals, and $r_{i}$ and $%
c_{j}$ the row and column averages; the values of $r_{i}$ and $c_{j}$ for
rows and columns with zero weight will not matter).

For every concave function $F:\mathbb{R}^{m}\rightarrow \mathbb{R}$ (it
suffices for $F$ to be defined on the compact convex set $\mathrm{conv}%
\{x_{ij}:i\in I,j\in J\}$) let%
\begin{eqnarray*}
E_{W}(F) &%
%TCIMACRO{\TeXButton{:=}{{\;:=\;}}}%
%BeginExpansion
{\;:=\;}%
%EndExpansion
&\sum_{i\in I}\sum_{j\in J}w_{ij}F(x_{ij}) \\
R_{W}(F) &%
%TCIMACRO{\TeXButton{:=}{{\;:=\;}}}%
%BeginExpansion
{\;:=\;}%
%EndExpansion
&\sum_{i\in I}w_{i\cdot }F(r_{i}) \\
C_{W}(F) &%
%TCIMACRO{\TeXButton{:=}{{\;:=\;}}}%
%BeginExpansion
{\;:=\;}%
%EndExpansion
&\sum_{j\in J}w_{\cdot j}F(c_{j})
\end{eqnarray*}%
(these are the overall average, average by rows, and average by columns,
respectively). We are interested in the inequalities $C\leq R$; more
precisely, when does the inequality $C(Q)\leq R(Q)$ for a quadratic concave
function $Q$ imply $C(F)\leq R(F)$ for all concave functions $F$.

The matrix $X$ is $W$\emph{-column-constant}\footnote{%
Interpreting $W$ as a probability measure, this means that $X$ is $W$%
-a.s.~column-constant.} if the restriction of $X$ to the support of $W$ has
constant columns, i.e., $x_{ij}=x_{i^{\prime }j}$ whenever $w_{ij}>0$ and $%
w_{i^{\prime }j}>0$; the matrix $X$ is \emph{column-constant }if it has
constant columns, i.e., $x_{ij}=x_{i^{\prime }j}$ for every $i,i^{\prime
}\in I$ and $j\in J$ (thus, $X$ is column-constant if and only if it is $W$%
-column-constant for every $W$). Similarly for $W$\emph{-row-constant} ($%
x_{ij}=x_{ij^{\prime }}$ whenever $w_{ij}>0$ and $w_{ij^{\prime }}>0$) and 
\emph{row-constant} ($x_{ij}=x_{ij^{\prime }}$ for every $i\in I$ and $%
j,j^{\prime }\in J$). For instance, if $W$ is a diagonal matrix then every $%
X $ is both $W$-column-constant and $W$-row-constant.

\begin{proposition}
\label{p:0}Let $W$ be a weight matrix. Then:

\begin{description}
\item[(a)] $C_{W}(F)\geq E_{W}(F)$ and $R_{W}(F)\geq E_{W}(F)$ for every
concave $F.$

\item[(b1)] If $X$ is $W$-column-constant then $C_{W}(F)=E_{W}(F)$ for every
concave $F$.

\item[(b2)] If $X$ is $W$-row-constant then $R_{W}(F)=E_{W}(F)$ for every
concave $F.$

\item[(c1)] $C_{W}(G)=E_{W}(G)$ for some \emph{strictly} concave function $G$
if and only if $X$ is $W$-column-constant (and then $C_{W}(F)=E_{W}(F)$ for
every concave $F$ by (b1)).

\item[(c2)] $R_{W}(G)=E_{W}(G)$ for some \emph{strictly} concave function $G$
if and only if $X$ is $W$-row-constant (and then $R_{W}(F)=E_{W}(F)$ for
every concave $F$ by (b2)).
\end{description}
\end{proposition}

\begin{proof}
(a) The concavity of $F$ yields%
\begin{equation}
F(c_{j})\geq \sum_{i}\frac{w_{ij}}{w_{\cdot j}}F(x_{ij})  \label{eq:concave}
\end{equation}%
for every $j$; multiplying by $w_{\cdot j}$ and summing over $j$ then gives%
\begin{equation}
C_{W}(F)=\sum_{j}w_{\cdot j}F(c_{j})\geq
\sum_{j}\sum_{i}w_{ij}F(x_{ij})=E_{W}(F).  \label{eq:R>=E}
\end{equation}

(b) If $X$ is $W$-column-constant then $x_{ij}=c_{j}$ for every $i,j$ with $%
w_{ij}>0$, and so we get equality in (\ref{eq:concave}), and thus in (\ref%
{eq:R>=E}). Similarly for (b2).

(c) Equality in (\ref{eq:R>=E}), and thus in (\ref{eq:concave}) for every $j$
with $w_{\cdot j}>0$, holds for a strictly concave $G$ if and only if all
the entries $x_{ij}$ in column $j$ that have a positive weight $w_{ij}>0$
must be equal; that is, $X$ is $W$-column-constant. Similarly for (c2).
\end{proof}

\bigskip

A useful strictly concave function $F$ is the quadratic $Q(z)=-\left\Vert
z\right\Vert ^{2}$; in fact, any other concave quadratic function works just
as well.\footnote{\label{ft:Q}Any inequality between $C_{W}(F)$ and $%
R_{W}(F) $ that holds for $F=Q$ holds also for every other concave quadratic
function $F$. Indeed, such an $F$ can be expressed as $F=\lambda Q+G$ for $%
\lambda >0$ and $G$ affine; since $C_{W}$ and $R_{W}$ are linear in $F$,
and, for affine $G$, we have $C_{W}(G)=E_{W}(G)=R_{W}(G)$, it follows that $%
C_{W}(\lambda Q+G)-R_{W}(\lambda Q+G)=\lambda (C_{W}(Q)-R_{W}(Q))$.}

We will say that a matrix $X$ is \emph{non-degenerate} if it has more than 2
distinct entries.

\begin{theorem}
\label{th:1}Let $X$ be a non-degenerate matrix. Then the following four
statements are equivalent.

\begin{description}
\item[(i)] For every weight matrix $W$, if $C_{W}(Q)\leq R_{W}(Q)$ then $%
C_{W}(F)\leq R_{W}(F)$ for every concave function $F$.

\item[(ii)] For every weight matrix $W$, if $C_{W}(Q)\leq R_{W}(Q)$ then $%
C_{W}(Q)=E_{W}(Q)$.

\item[(iii)] For every weight matrix $W$, if $C_{W}(Q)\leq R_{W}(Q)$ then $%
C_{W}(F)=E_{W}(F)\leq R_{W}(F)$ for every concave function $F$.

\item[(iv)] The matrix $X$ is column-constant or row-constant.
\end{description}

\noindent Moreover, when $X$ is column-constant, $C_{W}(F)=E_{W}(F)\leq
R_{W}(F)$ (for every concave $F$), and thus $C_{W}(Q)\leq R_{W}(Q)$, hold
for every $W$. When $X$ is row-constant, $C_{W}(Q)\leq R_{W}(Q)$ holds if
and only if $W$ is such that $X$ is also $W$-column-constant, and then $%
C_{W}(Q)=R_{W}(Q)$ and $C_{W}(F)=E_{W}(F)=R_{W}(F)$ (for every concave $F$).
\end{theorem}

\begin{proof}
The substantial part of the proof consists of showing that (i) implies (iv);
the rest then follows easily by using Proposition \ref{p:0}.

\textbf{(i) implies (iv). }Assume (i), i.e., for every $W$ for which $%
C_{W}(Q)\leq R_{W}(Q)$, it holds that $C_{W}(F)\leq R_{W}(F)$ for every
concave function $F$.

$\bullet $ \emph{Claim 1}$.$ For every $W$ such that $C_{W}(Q)\leq R_{W}(Q)$%
, the row averages $r_{i}$ must lie in the convex hull of the column
averages $c_{j}$; i.e., letting $I_{0}:=\{i\in I:w_{i\cdot }>0\}$ and $%
J_{0}:=\{j\in J:w_{\cdot j}>0\}$, we have $r_{i}\in \mathrm{conv}%
\{c_{j}:j\in J_{0}\}$ for every $i\in I_{0}$.

\emph{Proof. }Let $\Gamma :=\mathrm{conv}\{c_{j}:j\in J_{0}\}$ and assume
that $r_{i}\notin \Gamma $ for some $i$ with $w_{i\cdot }>0.$ Since $\Gamma $
is a closed convex set, by the Separating Hyperplane Theorem there are $v\in 
\mathbb{R}^{m}$ and $\alpha \in \mathbb{R}$ such that $v\cdot r_{i}>\alpha
\geq v\cdot c_{j}$ for every $j\in J_{0}$. For every $n$ let $F_{n}$ be the
following function: $F_{n}(z):=Q(z)-n(\max \{v\cdot z-\alpha ,0\})^{2}$; the
function $F_{n}$ is concave (as the sum of two concave functions) and
continuously differentiable.\footnote{\label{ft:smooth}The continuous
differentiability is needed for our application below to calibeating.} Now $%
C_{W}(F_{n})=C_{W}(Q)$ for all $n$ (because $F_{n}(c_{j})=Q(c_{j})$ for all $%
j$ with $w_{\cdot j}>0$), whereas $R_{W}(F_{n})\rightarrow -\infty $ as $%
n\rightarrow \infty $ (because $v\cdot r_{i}-\alpha >0$ and so $%
F_{n}(r_{i})\rightarrow -\infty $); therefore, for $n$ large enough the
inequality $C_{W}(F_{n})\leq R_{W}(F_{n})$ fails---contradicting (i).~$%
\square $

$\bullet $ \emph{Claim 2. }Every $2\times 2$ submatrix of $X$ has at most $2$
distinct entries.

\emph{Proof}. Let $X^{0}=\left[ 
\begin{array}{cc}
a & b \\ 
d & \cdot%
\end{array}%
\right] $ be a $2\times 2$ submatrix of $X$ with $a,b,d$ distinct (the
fourth entry will get weight $0$ and so will not matter). We distinguish two
cases, according to whether or not $d$ lies in the open interval $(a,b)$
(i.e., on the straight line through $a$ and $b$, between $a$ and $b$).

Case A: $d\notin (a,b).$ Let $W=\left[ 
\begin{array}{cc}
(1-\varepsilon )/2 & (1-\varepsilon )/2 \\ 
\varepsilon & 0%
\end{array}%
\right] $ for $\varepsilon >0$. As $\varepsilon \rightarrow 0$ we get $%
R_{W}(Q)\rightarrow -\left\Vert (a+b)/2\right\Vert ^{2}$ and $%
C_{W}(Q)\rightarrow -\left\Vert a\right\Vert ^{2}/2-\left\Vert b\right\Vert
^{2}/2,$ and so $a\neq b$ yields $C_{W}(Q)<R_{W}(Q)$ for small enough $%
\varepsilon >0$; but $c_{1}\rightarrow a,$ $c_{2}=b$, and $r_{2}=d$, and so $%
r_{2}\notin \mathrm{conv}\{c_{1},c_{2}\}$ for small enough $\varepsilon >0$,
a contradiction to Claim 1.

Case B: $d\in (a,b)$. Restricting to the $1$-dimensional space (the line)
that contains $a,b,$ and $d$ and using the coordinate system $\lambda
a+(1-\lambda )b\mapsto \lambda $ yields the matrix\footnote{%
The change of coordinates has no effect on the sign of $R(Q)-C(Q)$; it
amounts to using the quadratic $Q_{1}(z)=-\left\Vert z-b\right\Vert
^{2}/\left\Vert a-b\right\Vert ^{2}$ instead of $Q(z)=-\left\Vert
z\right\Vert ^{2}$ (cf. footnote \ref{ft:Q}).} $\left[ 
\begin{array}{cc}
1 & 0 \\ 
\delta & \cdot%
\end{array}%
\right] $ with $0<\delta <1.$ Let $W=\frac{1}{3(1+\delta )}\left[ 
\begin{array}{cc}
1+2\delta & 1-\delta \\ 
1+2\delta & 0%
\end{array}%
\right] ,$ then a straightforward computation yields $R_{W}(Q)-C_{W}(Q)=%
\delta (1-\delta )(1+2\delta )/(6(2+\delta ))>0$, but $c_{1}=(1+\delta )/2,$ 
$c_{2}=0$, and $r_{1}=(1+2\delta )/(2+\delta )>c_{1}>c_{2}$, and so $%
r_{1}\notin \mathrm{conv}\{c_{1},c_{2}\}$, a contradiction to Claim 1.

This completes the proof of Claim 2.~$\square $

$\bullet $ \emph{Claim 3.} $X$ is column-constant or row-constant (i.e.,
(iv)).

\emph{Proof}. Assume that $X$ is not column-constant; let $a\neq b$ be two
distinct entries in, say, column 1. Then Claim 2 implies that every row that
has $a$ or $b$ in column 1 must contain only $a$ and $b$. If column 1 were
to contain only $a$ and $b,$ then the entire matrix would contain only $a$
and $b$, a contradiction to the non-degeneracy of $X$. Therefore column 1
must contain another distinct entry $d$. But then each row with $a$ in
column 1 cannot contain $b$ (again by Claim 2, using the two rows with $a$
and $d$ in column 1), and so the row is a constant $a$ row; similarly, each
row with $b$ in column 1 cannot contain $a$, and so it is a constant $b$
row. Carrying out the same argument, but now with the pair of distinct
entries $a$ and $d$, shows that every row with $d$ in column 1 must be a
constant $d$ row; this holds for every $d$, and so all rows are constant
rows.~$\square $

\textbf{(iv) implies (iii) and the \textquotedblleft
moreover\textquotedblright\ statement.}

If $X$ is column-constant then $C_{W}(F)=E_{W}(F)\leq R_{W}(F)$ for every $W$
and every concave $F$ by Proposition \ref{p:0} (a) and (b1).

If $X$ is row-constant then $C_{W}(Q)\geq E_{W}(Q)=R_{W}(Q)$ for all $W$ by
Proposition \ref{p:0} (a) and (b2), and so if $W$ is such that $C_{W}(Q)\leq
R_{W}(Q)$ then we have equality $C_{W}(Q)=E_{W}(Q)=R_{W}(Q)$; since $Q$ is
strictly concave, by Proposition \ref{p:0} (b1), $X$ is (also) $W$%
-column-constant, and then $C_{W}(F)=E_{W}(F)=R_{W}(F)$ for every concave $F$
by Proposition \ref{p:0} (b1) and (b2).

Altogether, we have shown (iii), and also the \textquotedblleft
moreover\textquotedblright\ statement. $\square $

\textbf{(iii) implies (ii).} Immediate, since $Q$ is concave.~$\square $

\textbf{(ii) implies (i).} Use Proposition \ref{p:0} (c1): since $Q$ is
strictly concave, $C_{W}(Q)=E_{W}(Q)$ implies that $X$ is $W$%
-column-constant and then $C_{W}(F)=E_{W}(F)\leq R_{W}(F)$ for every concave 
$F$.
\end{proof}

\bigskip

\noindent \textbf{Remarks. }\emph{(a) }The result continues to hold even if
we restrict the concave functions $F$ to have Lipschitz-continuous gradients
in the relevant domain (namely, the compact convex set that contains all $%
x_{ij}$); indeed, our construction uses only such functions $F_{n}$ (see
Footnote \ref{ft:smooth}; continuous differentiability on a compact set
implies that the gradient is Lipschitz-continuous there).

\emph{(b)} The result is false when $X$ has only two distinct entries. For
example, let $X=\left[ 
\begin{array}{cc}
1 & 0 \\ 
0 & 1%
\end{array}%
\right] $. For every weight matrix $W$ we have $r_{1}\geq c_{1}$ iff $%
w_{11}/(w_{11}+w_{12})\geq w_{11}/(w_{11}+w_{21})$ iff $w_{12}\leq w_{21}$
iff $w_{22}/(w_{22}+w_{12})\geq w_{22}/(w_{22}+w_{21})$ iff $c_{2}\geq r_{2}$%
, and so $r_{1},r_{2}$ are either both inside the interval $c_{1},c_{2}$, or
both strictly outside. Because the average of the $r_{i}$ equals the average
of the $c_{j}$, in the first case we have $C_{W}(F)\leq R_{W}(F)$ for every
concave $F$, and in the second case we have $C_{W}(F)>R_{W}(F)$ for all
strictly concave $F.$ Therefore, (i) holds. However, (ii) does not hold:
take for instance all weights to be $1/4$, then $C_{W}(Q)=R_{W}(Q)=-1/4,$
whereas $E_{W}(Q)=-1/2$.

\subsection{Application to Proper Calibeating\label{s-s-a:bxc calibeating}}

Let $\mathbf{a}$ be an action sequence, and $\mathbf{b}$ and $\mathbf{c}$
two forecasting sequences. Ignoring for simplicity the horizon and the $o(1)$
errors, $\mathbf{c}$ \emph{calibeats} $\mathbf{b}$ if $\mathcal{B}(\mathbf{c}%
)\leq \mathcal{R}(\mathbf{b})$, and $\mathbf{c}$ \emph{proper-calibeats} $%
\mathbf{b}$ if $\mathcal{B}^{L}(\mathbf{c})\leq \mathcal{R}^{L}(\mathbf{b})$
for every proper -Lipschitz scoring rule $L$ (i.e., $L$ in $\mathcal{L}_{1}$%
).

Let $B$ and $C$ be the finite sets of bins of $\mathbf{b}$ and $\mathbf{c}$,
respectively, and let $\lambda (b,c)$ be the relative frequency of the $(b,c)
$-bin, and $\bar{a}(b,c)$ the average action there. Then\footnote{%
Let $\lambda (b,\cdot ):=\sum_{c}\lambda (b,c)$ and $\lambda (\cdot
,c):=\sum_{b}\lambda (b,c)$ denote the marginal frequencies, and $\bar{a}%
(b,\cdot )$ and $\bar{a}(\cdot ,c)$ the action averages, for the $b$-bins
and the $c$-bins, respectively.}%
\begin{eqnarray*}
\mathcal{R}^{L}(\mathbf{b}) &=&\sum_{b\in B}\lambda (b,\cdot )H^{L}(\bar{a}%
(b,\cdot ))-\mathcal{H}^{L} \\
\mathcal{R}^{L}(\mathbf{c}) &=&\sum_{c\in C}\lambda (\cdot ,c)H^{L}(\bar{a}%
(\cdot ,c))-\mathcal{H}^{L}
\end{eqnarray*}%
where $H^{L}$ is the corresponding $L$-entropy function, which is a concave
function, and $\mathcal{H}^{L}$ is the average $L$-entropy of the actions.
For the standard quadratic scoring rule, $H(z)=Q(z)=-\left\Vert z\right\Vert
^{2}.$

For the result below, we fix the action averages $\bar{a}(b,c)$, while
allowing the relative frequencies $\lambda (b,c)$ to vary. Since this
changes the $c$-bin averages, namely, $\bar{a}(\cdot ,c)=\sum_{b\in
B}\lambda (b,c)\bar{a}(b,c)/\sum_{b\in B}\lambda (b,c)$, in order for the
forecasting sequence $\mathbf{c}\equiv \mathbf{c}_{\lambda }$ to be
calibrated we replace each label $c$ with the corresponding $\bar{a}(\cdot
,c)$. While bin labels do not matter for refinement scores, the above
relabeling may change the binning itself, when two different $c$ bins get
the same new label. If, say, $\bar{a}(.,c)=\bar{a}(\cdot ,c^{\prime })$,
then we need to merge the overall $c$ and $c^{\prime }$ bins, and also the $%
(b,c)$ and $(b,c^{\prime })$ bins for every $b\in B$. Using the notation $%
\mathcal{R}^{L}(B),$ $\mathcal{R}^{L}(C),$ and $\mathcal{R}^{L}(B\times C)$
for the original refinement scores (i.e., before the above merging), we have 
$\mathcal{R}^{L}(\mathbf{b})=\mathcal{R}^{L}(B)$ (because the merging has no
effect on the overall $b$-bins), $\mathcal{R}^{L}(\mathbf{c})=\mathcal{R}%
^{L}(C)$ (because we merge total $c$-bins that have the same action
average), and $\mathcal{R}^{L}(\mathbf{b}\times \mathbf{c})\geq \mathcal{R}%
^{L}(B\times C)$ (because the $B\times C$ binning can only be finer).

\begin{theorem}
\label{th:bxc}Assume that the matrix $(\bar{a}(b,c))_{b\in B,c\in C}$ of
action averages is non-degenerate. Then the following statements are
equivalent:

\begin{description}
\item[(i)] For every bin-frequency matrix $\lambda $ and corresponding
calibrated forecasting sequence $\mathbf{c}\equiv \mathbf{c}_{\lambda }$, if 
$\mathbf{c}$ calibeats $\mathbf{b}$ then $\mathbf{c}$ proper-calibeats $%
\mathbf{b}$.

\item[(ii)] For every bin-frequency matrix $\lambda $ and corresponding
calibrated forecasting sequence $\mathbf{c}\equiv \mathbf{c}_{\lambda }$, if 
$\mathbf{c}$ calibeats $\mathbf{b}$ then $\mathbf{c}$ calibeats $\mathbf{%
b\times c}$ (and thus $\mathbf{c}$ is proper-calibrated, and
proper-calibeats $\mathbf{b\times c}$).
\end{description}
\end{theorem}

\begin{proof}
We use Theorem \ref{th:1}. Let $X=(\bar{a}(b,c))_{b\in B,c\in C}$ be the
matrix of action averages, and $W=\lambda \equiv (\lambda (b,c))_{b\in
B,c\in C}$ the matrix of bin frequencies. Then $\mathcal{R}^{L}(\mathbf{b})=%
\mathcal{R}^{L}(B)=R_{W}(H^{L})-\mathcal{H}^{L}$, $\mathcal{R}^{L}(\mathbf{c}%
)=\mathcal{R}^{L}(C)=C_{W}(H^{L})-\mathcal{H}^{L}$, and $\mathcal{R}^{L}(%
\mathbf{b}\times \mathbf{c})\geq \mathcal{R}^{L}(B\times C)=E_{W}(H^{L})-%
\mathcal{H}^{L}$ (for every $L$). Statement (i) of Theorem \ref{th:1}
translates to: for every $\lambda $, if $\mathcal{R}(\mathbf{c})\leq 
\mathcal{R}(\mathbf{b})$ then $\mathcal{R}^{L}(\mathbf{c})\leq \mathcal{R}%
^{L}(\mathbf{b})$ for every\footnote{%
See Remark (a) above: the restriction to Lipschitz $L$ does not matter for
the result.} $L$ in $\mathcal{L}$; this is statement (i) here. Statement
(ii) of Theorem \ref{th:1} translates to: for every $\lambda $, if $\mathcal{%
R}(\mathbf{c})\leq \mathcal{R}(\mathbf{b})$ then $\mathcal{R}(\mathbf{c})=%
\mathcal{R}(B\times C)\leq \mathcal{R}(\mathbf{b}\times \mathbf{c})$; this
is statement (ii) here (use (J1) in the Remark following Theorem \ref%
{th:beat-by-calib S}).
\end{proof}

\bigskip 

\textbf{[TO BE REWRITTEN] }One way to interpret this is as follows: if, for
some bin frequencies, a calibrated $\mathbf{c}$ calibeating $\mathbf{b}$
entails that it proper-calibeats $\mathbf{b}$, but it does not calibeat the
joint $\mathbf{b}\times \mathbf{c}$, then one can change the frequencies
such that the resulting calibrated $\mathbf{c}^{\prime }$ calibeats $\mathbf{%
b}$ but does \emph{not} proper-calibeat $\mathbf{b}$.

\begin{thebibliography}{99}
\bibitem{} \textbf{[NEEDS\ UPDATING]}

\bibitem{} Azoury, K. S. and M. K. Warmuth (2001), \textquotedblleft
Relative Loss Bounds for On-Line Density Estimation with the Exponential
Family of Distributions,\textquotedblright\ \emph{Machine Learning} 43,
211--246.

\bibitem{} Blackwell, D. (1956), \textquotedblleft An Analog of the Minimax
Theorem for Vector Payoffs,\textquotedblright\ \emph{Pacific Journal of
Mathematics} 6, 1--8.

\bibitem{} Brier, G. W. (1950), \textquotedblleft Verification of Forecasts
Expressed in Terms of Probability,\textquotedblright\ \emph{Monthly Weather
Review} 78, 1--3.

\bibitem{} Brouwer, L. E. J. (1912), \textquotedblleft \"{U}ber Abbildung
von Mannigfaltigkeiten,\textquotedblright\ \emph{Mathematische Annalen} 71,
97--115.

\bibitem{} Cesa-Bianchi, N. and G. Lugosi (2006), \emph{Prediction,
Learning, and Games}, Cambridge University Press.

\bibitem{} Chen, Y., Z. Huang, M. I. Jordan, and H. Luo, \textquotedblleft
Calibeating Made Simple,\textquotedblright\ (2026), \texttt{arXiv}%
:2603.22167.

\bibitem{} Dawid, A. (1982), \textquotedblleft The Well-Calibrated
Bayesian,\textquotedblright\ \emph{Journal of the American Statistical
Association} 77, 605--613.

\bibitem{} Forster, J. (1999), \textquotedblleft On Relative Loss Bounds in
Generalized Linear Regression,\textquotedblright\ in \emph{12th
International Symposium on Fundamentals of Computation Theory (FCT '99)},
269--280.

\bibitem{} Foster, D. P. (1991), \textquotedblleft Prediction in the Worst
Case,\textquotedblright\ \emph{The Annals of Statistics} 19, 1084--1090.

\bibitem{} Foster, D. P. (1999), \textquotedblleft A Proof of Calibration
via Blackwell's Approachability Theorem,\textquotedblright\ \emph{Games and
Economic Behavior }29, 73--78.

\bibitem{} Foster, D. P. and S. Hart (2018), \textquotedblleft Smooth
Calibration, Leaky Forecasts, Finite Recall, and Nash
Dynamics,\textquotedblright\ \emph{Games and Economic Behavior }109,
271--293.

\bibitem{} Foster, D. P. and S. Hart (2021), \textquotedblleft Forecast
Hedging and Calibration,\textquotedblright\ \emph{Journal of Political
Economy} 129, 3447--3490. doi.org/10.1086/716559. \texttt{%
http://www.ma.huji.ac.il/hart/publ.html\#calib-int}

\bibitem{} Foster, D. P. and S. Hart (2023), \textquotedblleft\
`Calibeating': Beating Forecasters at Their Own Game,\textquotedblright\ 
\emph{Theoretical Economics }18, 1441--1474.

\bibitem{} Foster, D. P. and S. Hart (2026), \textquotedblleft Errata:\
`Calibeating': Beating Forecasters at Their Own Game,\textquotedblright\ 
\texttt{http://www.ma.huji.ac.il/hart/papers/calib-beat-errata.pdf}

\bibitem{} Foster, D. P. and R. V. Vohra (1998), \textquotedblleft
Asymptotic Calibration,\textquotedblright\ \emph{Biometrika} 85, 379--390.

\bibitem{} Hart, S. (1992), \textquotedblleft Games in Extensive and
Strategic Forms,\textquotedblright\ in \emph{Handbook of Game Theory, with
Economic Applications}, R. J. Aumann and S. Hart (editors), North-Holland,
Vol. 1, Chapter 2, 19--40.

\bibitem{} Hart, S. (2021), \textquotedblleft Calibrated Forecasts: The
Minimax Proof,\textquotedblright\ Center for Rationality DP-744, The Hebrew
University of Jerusalem. \texttt{http://www.ma.huji.ac.il/hart/publ.html%
\#calib-minmax}

\bibitem{} Johnson, N. L., A. W. Kemp, and S. Kotz (2005), \emph{Univariate
Discrete Distributions}, 3rd edition, Wiley.

\bibitem{} Jung, H. (1901), \textquotedblleft Ueber die kleinste Kugel, die
eine r\"{a}umliche Figur einschliesst,\textquotedblright\ \emph{Journal f%
\"{u}r die Reine und Angewandte Mathematik} 123, 241--257.

\bibitem{} Kuhn, H. W. (1953), \textquotedblleft Extensive Games and the
Problem of Information,\textquotedblright\ in \emph{Contributions to the
Theory of Games, Vol. II}, H. W. Kuhn and A. W. Tucker (editors), \emph{%
Annals of Mathematics Studies} 28, Princeton University Press, 193--216.

\bibitem{} Lee, D., Noarov, G., Pai, M., and Roth, A. (2022),
\textquotedblleft Online Minimax Multiobjective Optimization:
Multicalibeating and Other Applications,\textquotedblright\ in: \emph{%
Advances in Neural Information Processing Systems}, 29051--29063. \texttt{%
https://proceedings.neurips.cc/paper%
\_files/paper/2022/file/ba942323c447c9bbb9d4b638eadefab9-Paper-Conference.pdf%
}

\bibitem{} Lo\`{e}ve, M. (1978), \emph{Probability Theory, Vol. II}, 4th
edition, Springer.

\bibitem{} Murphy, A. H. (1972), \textquotedblleft Scalar and Vector
Partitions of the Probability Score. Part I: Two-State
Situation,\textquotedblright\ \emph{Journal of Applied Meteorology} 11,
273--282.

\bibitem{} Oakes, D. (1985), \textquotedblleft Self-Calibrating Priors Do
Not Exist,\textquotedblright\ \emph{Journal of the American Statistical
Association} 80, 339.

\bibitem{} Olszewski, W. (2015), \textquotedblleft Calibration and Expert
Testing,\textquotedblright\ in \emph{Handbook of Game Theory, Vol. 4}, H. P.
Young and S. Zamir (editors), Springer, 949--984.

\bibitem{} Olszewski, W. and A. Sandroni (2008), \textquotedblleft
Manipulability of Future-Independent Tests,\textquotedblright\ \emph{%
Econometrica} 76, 1437--1466.

\bibitem{} Sanders, F. (1963), \textquotedblleft On Subjective Probability
Forecasting,\textquotedblright\ \emph{Journal of Applied Meteorology} 2,
191--201.

\bibitem{} Sandroni, A. (2003), \textquotedblleft The Reproducible
Properties of Correct Forecasts,\textquotedblright\ \emph{International
Journal of Game Theory} 32, 151--159.

\bibitem{} Shmaya, E. (2008), \textquotedblleft Many Inspections are
Manipulable,\textquotedblright\ \emph{Theoretical Economics} 3, 367--382.

\bibitem{} Welford, B. P. (1962), \textquotedblleft Note on a Method for
Calculating Corrected Sums of Squares and Products,\textquotedblright\ \emph{%
Technometrics} 4, 419--420.

\bibitem{} Vovk, V. (2001), \textquotedblleft Competitive On-Line
Statistics,\textquotedblright\ \emph{International Statistical Review} 69,
213--248.

\bibitem{} von Neumann, J. (1928), \textquotedblleft Zur Theorie der
Gesellschaftsspiele, \textquotedblright\ \emph{Mathematische Annalen} 100,
295--320.
\end{thebibliography}

\end{document}
